% =====================================================================
%  Capítulo 13 · Epigenómica y regulación
% =====================================================================
\chapter{Epigenómica y regulación}\label{cap:13}

\epigrafe{La naturaleza combinatoria de las modificaciones de los extremos amino de las histonas revela así un ``código de histonas'' que amplía considerablemente el potencial informativo del código genético.}{\textcite{c13-jenuwein2001}}

\begin{objetivos}
\begin{itemize}
  \item Describir la organización de la cromatina (nucleosomas, marcas de histonas, estados de cromatina) y los grandes atlas epigenómicos (ENCODE, Roadmap Epigenomics).
  \item Explicar la geometría de las lecturas de ChIP-seq, estimar el tamaño de fragmento por correlación cruzada y llamar picos con el modelo de Poisson local de MACS; evaluar un experimento con FRiP, NSC/RSC, IDR y la lista negra.
  \item Interpretar la distribución de tamaños de fragmento y el enriquecimiento en TSS de un experimento de ATAC-seq.
  \item Deducir la química de la conversión por bisulfito, alinear lecturas convertidas y cuantificar la metilación con valores $\beta$ y $M$; detectar regiones diferencialmente metiladas con modelos beta-binomiales.
  \item Formular el descubrimiento de motivos como un problema de datos incompletos y resolverlo con el algoritmo EM (MEME) y con muestreo de Gibbs; evaluar motivos por enriquecimiento diferencial y centralidad.
\end{itemize}
\end{objetivos}

Todas las células de un organismo llevan, salvo excepciones, el mismo genoma. Una neurona y un linfocito leen, sin embargo, capítulos muy distintos de ese texto común: la primera expresa canales iónicos y receptores de neurotransmisores; el segundo, receptores de antígenos y citocinas. La diferencia no está en la secuencia sino en \emph{cómo se lee}: qué regiones están accesibles a la maquinaria de transcripción, qué factores de transcripción se unen a ellas, qué marcas químicas llevan las histonas que empaquetan el ADN y qué citosinas están metiladas. Ese conjunto de estados, heredable a través de las divisiones celulares pero reversible, es el \emph{epigenoma}, y la \emph{epigenómica} es el intento de medirlo a escala de genoma completo.

La secuenciación masiva convirtió la epigenómica en una disciplina cuantitativa: casi todas sus técnicas terminan en millones de lecturas cortas que se mapean contra la referencia (\cref{cap:07}) y cuyo patrón de acumulación hay que interpretar estadísticamente. En este capítulo recorreremos las tres preguntas centrales del área. Primero, \emph{¿dónde} se une una proteína o dónde está abierta la cromatina? Lo responden ChIP-seq y ATAC-seq, y la herramienta matemática será el modelo de Poisson para distinguir un pico de un fondo (\cref{sec:13-chip}). Segundo, \emph{¿qué} citosinas están metiladas y cuánto? Lo responde la secuenciación con bisulfito, y la herramienta será el modelo beta-binomial (\cref{sec:13-metil}). Tercero, \emph{¿por qué} se une un factor allí y no en otro lugar? Lo responde el descubrimiento de motivos, que retoma las matrices de peso posicional del \cref{cap:04} y las aprende de los datos con los algoritmos EM y de Gibbs (\cref{sec:13-motivos}).

% ---------------------------------------------------------------------
\section{ChIP-seq y ATAC-seq}\label{sec:13-chip}
% ---------------------------------------------------------------------
\enlacecolab{modulo-13-epigenomica/13.1_chipseq_atacseq.ipynb}{13.1 · ChIP-seq y ATAC-seq}

\begin{intuicion}[Huellas en la nieve]
Imagine que quiere saber dónde se detienen los ciervos en un bosque nevado, pero nunca los ve: solo encuentra huellas al amanecer. Si caminara al azar por el bosque, vería huellas dispersas por todas partes; pero alrededor de los comederos las huellas se amontonan, llegan desde el norte por un lado y se van hacia el sur por el otro. Para decidir si un amontonamiento es un comedero o una casualidad necesita saber cuántas huellas \emph{esperaría} en esa zona si los ciervos pasearan sin rumbo (el fondo), y ese fondo no es igual en todo el bosque: junto al arroyo siempre hay más tránsito. ChIP-seq es exactamente esto. Las lecturas son huellas de fragmentos de ADN que estuvieron unidos a una proteína; se amontonan alrededor del sitio de unión, llegan por la hebra $+$ desde un lado y por la hebra $-$ desde el otro, y la estadística consiste en compararlas con un fondo local bien estimado.
\end{intuicion}

\subsection{La cromatina: un genoma empaquetado}

El ADN de una célula humana mide unos dos metros y cabe en un núcleo de pocas micras. Lo consigue enrollándose alrededor de octámeros de histonas (dos copias de cada una de H2A, H2B, H3 y H4). Cada unidad, el \emph{nucleosoma}, envuelve unos 147 pares de bases en una superhélice de casi dos vueltas; la estructura cristalográfica de \textcite{c13-luger1997} mostró en detalle atómico cómo 146 pares de bases se organizan alrededor del octámero y cómo las colas amino-terminales de las histonas salen entre las vueltas del ADN para contactar con partículas vecinas. Entre un nucleosoma y el siguiente queda un segmento de ADN \emph{conector} (\ingles{linker}) de longitud variable, de modo que la repetición nucleosomal típica es de unos 180 a 200 pb. Esta cifra reaparecerá, como una huella dactilar, en los datos de ATAC-seq.

Las colas de las histonas son el soporte de un rico repertorio de modificaciones postraduccionales: acetilación, metilación (mono-, di- o tri-), fosforilación, ubiquitinación. Se nombran con una notación compacta: \texttt{H3K27ac} es la acetilación de la lisina 27 de la histona H3; \texttt{H3K4me3}, la trimetilación de la lisina 4. \textcite{c13-jenuwein2001} propusieron que la combinación de marcas forma un \emph{código de histonas} que, leído por proteínas con dominios específicos, determina transiciones entre cromatina activa y silenciada. La \cref{tab:13-marcas} resume las marcas más estudiadas.

\begin{table}[t]
\centering\small
\caption[Marcas de histonas frecuentes]{Marcas de histonas más usadas en los atlas epigenómicos, con su localización típica y la forma de su señal en ChIP-seq. Las marcas ``anchas'' cubren decenas de kilobases y requieren estrategias de llamada distintas de las de un factor de transcripción.}
\label{tab:13-marcas}
\begin{tabularx}{\linewidth}{@{}l>{\raggedright\arraybackslash}X>{\raggedright\arraybackslash}p{2.1cm}l@{}}
\toprule
\textbf{Marca} & \textbf{Dónde aparece} & \textbf{Asociación} & \textbf{Señal}\\
\midrule
H3K4me3 & promotores, alrededor del TSS & activa & estrecha\\
H3K4me1 & potenciadores (\ingles{enhancers}) activos o preparados & activa/latente & intermedia\\
H3K27ac & potenciadores y promotores activos & activa & intermedia\\
H3K36me3 & cuerpos de genes transcritos & elongación & ancha\\
H3K27me3 & genes reprimidos por Polycomb & represiva & ancha\\
H3K9me3 & heterocromatina constitutiva, repeticiones & represiva & ancha\\
\bottomrule
\end{tabularx}
\end{table}

Tres grandes consorcios han cartografiado estas marcas junto con la unión de factores de transcripción y la accesibilidad de la cromatina. El proyecto ENCODE (\ingles{Encyclopedia of DNA Elements}) publicó en 2012 un mapa integrado de transcripción, unión de factores, estructura de la cromatina y modificación de histonas en líneas celulares humanas, y afirmó poder asignar alguna ``función bioquímica'' al \SI{80}{\percent} del genoma \parencite{c13-encode2012}. La cifra provocó un intenso debate, porque actividad bioquímica no equivale a función biológica; la lección metodológica es que \emph{toda señal epigenómica es una medida, no una interpretación}. La tercera fase del proyecto amplió la enciclopedia a humano y ratón y organizó los datos en un registro de elementos candidatos \emph{cis}-reguladores \parencite{c13-encode2020}. En paralelo, el consorcio Roadmap Epigenomics generó 111 epigenomas de referencia de tejidos y tipos celulares primarios e integró cinco marcas centrales (H3K4me3, H3K4me1, H3K36me3, H3K27me3 y H3K9me3) en un modelo de \emph{estados de cromatina} con el que anotar cada segmento del genoma como promotor activo, potenciador, región transcrita, región reprimida por Polycomb o heterocromatina \parencite{c13-kundaje2015}.

\subsection{El experimento de ChIP-seq}

La inmunoprecipitación de cromatina seguida de secuenciación (ChIP-seq) mide dónde está una proteína en el genoma de una población de células. \textcite{c13-johnson2007} la usaron para mapear la unión del represor NRSF (también llamado REST) a 1946 sitios del genoma humano, con una resolución de unas $\pm50$ pb, lo bastante fina para descubrir variantes no canónicas de su motivo. El protocolo tiene cinco pasos (\cref{fig:13-chip}):

\begin{enumerate}
  \item \textbf{Entrecruzamiento}. El formaldehído forma enlaces covalentes entre proteínas y ADN cercanos, ``congelando'' las interacciones del momento.
  \item \textbf{Fragmentación}. La cromatina se rompe por sonicación (o con nucleasa) en fragmentos de unos 100 a 300 pb.
  \item \textbf{Inmunoprecipitación}. Un anticuerpo específico contra la proteína (o contra una marca de histona) captura los fragmentos que la llevan; el resto se lava.
  \item \textbf{Reversión y secuenciación}. Se revierte el entrecruzamiento, se purifica el ADN, se prepara una biblioteca y se secuencia, casi siempre solo por un extremo de cada fragmento o en pares cortos.
  \item \textbf{Mapeo}. Las lecturas se alinean a la referencia (\cref{cap:07}) y se buscan regiones con más lecturas de las esperadas: los \emph{picos}.
\end{enumerate}

Siempre se acompaña de un \emph{control}: normalmente el \ingles{input}, ADN fragmentado de la misma muestra pero sin inmunoprecipitar, que registra los sesgos que nada tienen que ver con la proteína (regiones más fáciles de sonicar, amplificaciones del número de copias, zonas de mapeo problemático). \textcite{c13-park2009} ofrece una revisión clásica de las ventajas de ChIP-seq frente a su predecesor en micromatrices (ChIP-chip) y de los retos de diseño y análisis.

\begin{figure}[t]
\centering
\begin{tikzpicture}[font=\sffamily\scriptsize]
  \tikzset{nuc13/.style={circle,fill=gris!55,draw=gris,minimum size=4.2mm,inner sep=0pt},
           fc13/.style={fill=naranja,draw=naranja!70!black,rounded corners=1.5pt,minimum width=3.4mm,minimum height=2.6mm,inner sep=0pt}}
  \foreach \px/\tit in {0/{1. Entrecruzar},2.62/{2. Fragmentar},5.24/{3. Capturar (IP)},7.86/{4. Secuenciar},10.48/{5. Mapear}}{
    \draw[rejilla,rounded corners=3pt,fill=papel!35] (\px,0) rectangle (\px+2.45,3.3);
    \node[font=\sffamily\bfseries\scriptsize,text=tinta] at (\px+1.225,3.08) {\tit};}
  % 1: cromatina con nucleosomas y factor
  \draw[azul,thick] (0.15,1.5) .. controls (0.5,1.9) and (0.8,1.1) .. (1.2,1.5) .. controls (1.6,1.9) and (1.9,1.1) .. (2.3,1.5);
  \node[nuc13] at (0.45,1.62) {}; \node[nuc13] at (1.95,1.38) {};
  \node[fc13] at (1.2,1.72) {};
  \foreach \x in {1.12,1.28}{\node[text=rojooscuro,font=\sffamily\tiny\bfseries] at (\x,1.5) {x};}
  \node[etiqueta,align=center] at (1.225,0.55) {formaldehído\\fija proteína--ADN};
  % 2: fragmentos
  \draw[azul,thick] (2.8,2.2) -- (3.5,2.2); \node[nuc13] at (3.15,2.33) {};
  \draw[azul,thick] (3.7,1.6) -- (4.5,1.6); \node[fc13] at (4.1,1.76) {};
  \draw[azul,thick] (2.9,1.0) -- (3.6,1.0);
  \draw[azul,thick] (3.9,2.5) -- (4.8,2.5);
  \node[etiqueta] at (3.845,0.45) {sonicación};
  % 3: IP con anticuerpo y esfera
  \draw[azul,thick] (5.65,1.25) -- (6.55,1.25); \node[fc13] (f3) at (6.1,1.41) {};
  \draw[violeta,very thick] (6.1,1.55) -- (6.1,1.95) -- (5.9,2.2); \draw[violeta,very thick] (6.1,1.95) -- (6.3,2.2);
  \shade[ball color=amarillo!80] (6.1,2.55) circle (0.28);
  \node[etiqueta,text=violeta] at (6.85,1.95) {anticuerpo};
  \node[etiqueta] at (6.95,2.65) {esfera};
  \draw[azul!35,thick] (5.45,0.75) -- (6.1,0.75); \draw[azul!35,thick] (6.35,0.55) -- (7.1,0.55);
  \node[etiqueta,text=gris] at (6.465,0.25) {se lava lo no unido};
  % 4: secuenciación de ambos extremos
  \draw[azul,thick] (8.2,1.7) -- (9.9,1.7);
  \fill[naranja!70] (8.1,1.62) rectangle (8.2,1.78); \fill[naranja!70] (9.9,1.62) rectangle (10.0,1.78);
  \draw[-{Stealth[length=2mm]},azul,very thick] (8.2,2.1) -- (8.85,2.1);
  \draw[-{Stealth[length=2mm]},rojo,very thick] (9.9,1.3) -- (9.25,1.3);
  \node[etiqueta,text=azul] at (8.55,2.4) {lectura $+$};
  \node[etiqueta,text=rojooscuro] at (9.55,0.95) {lectura $-$};
  \node[etiqueta] at (9.085,0.45) {extremos 5'};
  % 5: mapeo
  \draw[tinta2] (10.6,0.8) -- (12.8,0.8);
  \draw[naranja,thick,dashed] (11.7,0.7) -- (11.7,2.4);
  \foreach \y/\x in {1.05/10.95,1.25/11.05,1.45/10.85,1.65/11.15,1.85/11.0,2.05/10.9}{
    \draw[-{Stealth[length=1.4mm]},azul,thick] (\x,\y) -- (\x+0.4,\y);}
  \foreach \y/\x in {1.15/12.45,1.35/12.3,1.55/12.5,1.75/12.35,1.95/12.4,2.15/12.25}{
    \draw[-{Stealth[length=1.4mm]},rojo,thick] (\x,\y) -- (\x-0.4,\y);}
  \node[etiqueta,text=naranja!80!black,fill=papel!35,inner sep=1pt] at (11.7,2.55) {sitio};
  \node[etiqueta] at (11.705,0.45) {genoma};
\end{tikzpicture}
\caption[Flujo de un experimento de ChIP-seq]{Flujo de un experimento de ChIP-seq. (1) El formaldehído fija el factor de transcripción (naranja) a su sitio; los nucleosomas (gris) también quedan fijados. (2) La cromatina se fragmenta. (3) Un anticuerpo acoplado a una esfera magnética captura los fragmentos con el factor; los demás se descartan. (4) Tras purificar el ADN se secuencia un extremo de cada fragmento: si se lee la hebra superior, el extremo 5' está a la izquierda del sitio; si se lee la inferior, a la derecha. (5) Al mapear, las lecturas $+$ (azul) se acumulan a la izquierda del sitio y las $-$ (rojo) a la derecha: la firma geométrica de un sitio de unión verdadero.}
\label{fig:13-chip}
\end{figure}

\subsection{La geometría de las lecturas: dos hebras y un desplazamiento}

El detalle clave del paso 4 es que la secuenciación lee solo las primeras decenas de bases de cada fragmento, desde su extremo 5'. Si el fragmento tiene longitud $\ell$ y está centrado en el sitio de unión $\mu$, la lectura de la hebra $+$ empieza en $\mu-\ell/2$ y la de la hebra $-$ en $\mu+\ell/2$. Con muchos fragmentos, las lecturas de cada hebra forman una ``montaña'' desplazada del sitio.

\begin{teorema}{Perfiles esperados por hebra}{13-hebras}
Sea $c$ la posición del centro de un fragmento, con densidad $\varphi(c)$ centrada en el sitio $\mu$, y sea $\ell$ su longitud, independiente de $c$, con densidad $g(\ell)$. Las densidades esperadas de extremos 5' en cada hebra son
\begin{equation}\label{eq:13-hebras}
f_{+}(x) = \int g(\ell)\,\varphi\!\left(x+\tfrac{\ell}{2}\right)d\ell,\qquad
f_{-}(x) = \int g(\ell)\,\varphi\!\left(x-\tfrac{\ell}{2}\right)d\ell .
\end{equation}
Si $\varphi$ es simétrica alrededor de $\mu$, $f_+$ y $f_-$ son imágenes especulares una de otra y sus centros de masa están en $\mu\mp\bar\ell/2$; la distancia entre ambos es la longitud media del fragmento, $d=\bar\ell$.
\end{teorema}

\begin{simbolos}
\simbolo{$\mu$}{Posición del sitio de unión en el genoma.}
\simbolo{$\varphi(c)$}{Densidad de la posición del centro del fragmento; su anchura refleja la imprecisión de la fragmentación.}
\simbolo{$g(\ell),\ \bar\ell$}{Densidad y media de la longitud de los fragmentos de la biblioteca.}
\simbolo{$f_\pm(x)$}{Densidad esperada de extremos 5' de lecturas en la hebra $+$ o $-$ en la posición $x$.}
\simbolo{$d$}{Desplazamiento entre las dos montañas; se estima de los datos y se usa para ``centrar'' las lecturas.}
\end{simbolos}

La demostración es un cambio de variable: una lectura $+$ cae en $x$ si el fragmento tiene centro $c=x+\ell/2$, de modo que basta integrar sobre las longitudes posibles. El resultado tiene una consecuencia práctica inmediata. Si desplazamos cada lectura $+$ una distancia $d/2$ hacia la derecha y cada lectura $-$ una distancia $d/2$ hacia la izquierda, las dos montañas se superponen en el sitio verdadero y la resolución mejora. Alternativamente, se puede \emph{extender} cada lectura hasta la longitud $d$ en dirección 3', reconstruyendo el fragmento completo; la suma de fragmentos reconstruidos se llama \ingles{pileup}. MACS estima $d$ tomando las regiones más enriquecidas, separando sus lecturas por hebra y midiendo la distancia entre las modas de ambas \parencite{c13-zhang2008}. La \cref{fig:13-hebras}\,(a) muestra el resultado con fragmentos simulados de longitud media 200 pb: las modas están en $-100$ y $+99$, y la suma desplazada queda centrada en el sitio.

\subsection{Correlación cruzada entre hebras}

Una estimación de $d$ que no requiere elegir regiones de antemano usa \emph{todo} el genoma. Sean $n_+(x)$ y $n_-(x)$ los conteos de extremos 5' en cada posición. Si los sitios de unión dominan la señal, $n_-$ se parece a $n_+$ desplazado $d$ posiciones, y la correlación entre ambos perfiles será máxima en ese desplazamiento \parencite{c13-kharchenko2008}.

\begin{definicion}{Correlación cruzada entre hebras}{13-xcor}
Para cada desplazamiento $s\ge0$,
\begin{equation}\label{eq:13-xcor}
r(s) = \operatorname{corr}\bigl(n_{+}(x),\ n_{-}(x+s)\bigr)
= \frac{\sum_x \bigl(n_+(x)-\bar n_+\bigr)\bigl(n_-(x+s)-\bar n_-\bigr)}{(G-s)\,\sigma_+\,\sigma_-},
\end{equation}
donde la suma recorre las $G-s$ posiciones válidas del genoma. El \emph{tamaño de fragmento estimado} es $\hat d=\argmax_{s} r(s)$, excluyendo desplazamientos cercanos a la longitud de lectura.
\end{definicion}

\begin{simbolos}
\simbolo{$n_\pm(x)$}{Número de lecturas cuyo extremo 5' cae en $x$ en cada hebra.}
\simbolo{$\bar n_\pm,\ \sigma_\pm$}{Media y desviación estándar de esos conteos a lo largo del genoma.}
\simbolo{$G$}{Longitud del genoma (o del cromosoma analizado).}
\simbolo{$s$}{Desplazamiento aplicado a la hebra $-$.}
\end{simbolos}

La curva $r(s)$ tiene casi siempre \emph{dos} máximos. El primero, en $s=\hat d$, es la señal biológica. El segundo, en $s\approx L$ (la longitud de lectura), es el \emph{pico fantasma} (\ingles{phantom peak}), un artefacto de la mapabilidad: en regiones repetidas, donde las lecturas no pueden mapearse de forma única, desaparecen a la vez las lecturas $+$ que empiezan en $x$ y las lecturas $-$ cuyo extremo 5' está en $x+L-1$, lo que correlaciona ambas hebras a esa distancia aunque no haya proteína alguna. Las guías de ENCODE y modENCODE \parencite{c13-landt2012} convirtieron la curva en dos índices de calidad:
\begin{equation}\label{eq:13-nsc}
\mathrm{NSC} = \frac{r(\hat d)}{\min_s r(s)},\qquad
\mathrm{RSC} = \frac{r(\hat d)-\min_s r(s)}{r(L)-\min_s r(s)} .
\end{equation}
\begin{simbolos}
\simbolo{NSC}{Coeficiente de hebras normalizado: cuánto sobresale el pico de fragmento sobre el fondo de la curva.}
\simbolo{RSC}{Coeficiente de hebras relativo: tamaño del pico de fragmento comparado con el del pico fantasma.}
\simbolo{$L$}{Longitud de las lecturas.}
\end{simbolos}
Un RSC mayor que 1 indica que la señal biológica domina sobre el artefacto de mapeo; según esas guías, valores de NSC por debajo de 1,05 y de RSC por debajo de 0,8 delatan un enriquecimiento pobre. En la simulación de la \cref{fig:13-hebras}\,(b), con lecturas de 36 pb, la muestra buena alcanza su máximo en $\hat d=176$ pb con $\mathrm{NSC}=2{,}70$ y $\mathrm{RSC}=1{,}76$; en la muestra pobre, con siete veces menos lecturas procedentes de sitios reales, el pico de fragmento casi desaparece y el máximo fuera de la zona del fantasma cae en un desplazamiento de 104 pb sin significado biológico, con $\mathrm{RSC}=0{,}36$.

\begin{figure}[t]
\centering
\begin{tikzpicture}
\begin{groupplot}[group style={group size=2 by 1, horizontal sep=1.35cm},
   curso, height=5.3cm]
\nextgroupplot[width=6.4cm, title={(a) extremos 5' por hebra},
   xlabel={posición relativa al sitio (pb)}, ylabel={densidad (u.\,a.)},
   xmin=-350, xmax=350, ymin=0, ymax=2.1, xtick={-300,-200,-100,0,100,200,300}]
  \addplot[azul, very thick] table[x=x,y=plus] {figuras/cap13/hebras.dat};
  \addplot[rojo, very thick] table[x=x,y=minus] {figuras/cap13/hebras.dat};
  \addplot[tinta2, thick, dashed] table[x=x,y=comb] {figuras/cap13/hebras.dat};
  \draw[{Stealth}-{Stealth},naranja!80!black,thick] (axis cs:-100,1.8) -- node[above,etiqueta,text=naranja!70!black]{$d\approx200$ pb} (axis cs:100,1.8);
\nextgroupplot[width=6.4cm, title={(b) correlación cruzada},
   xlabel={desplazamiento $s$ (pb)}, ylabel={$r(s)\times 100$},
   xmin=0, xmax=400, ymin=1.5, ymax=7.4, xtick={0,100,200,300,400},
   legend style={at={(0.98,0.98)},anchor=north east}]
  \addplot[azul, very thick] table[x=s,y=bueno] {figuras/cap13/xcor.dat};
  \addlegendentry{buena}
  \addplot[naranja, very thick] table[x=s,y=pobre] {figuras/cap13/xcor.dat};
  \addlegendentry{pobre}
  \draw[tinta2,dashed] (axis cs:36,1.5) -- (axis cs:36,5.2);
  \node[etiqueta,anchor=south west] at (axis cs:10,5.2) {fantasma ($L$)};
  \draw[azul!70!black,dashed] (axis cs:176,1.5) -- (axis cs:176,6.7);
  \node[etiqueta,anchor=south,text=azulprofundo] at (axis cs:176,6.7) {$\hat d=176$};
\end{groupplot}
\end{tikzpicture}
\caption[Geometría de las lecturas de ChIP-seq]{(a) Densidad de extremos 5' de \num{4000} fragmentos simulados alrededor de un sitio de unión (longitud media 200 pb). Las lecturas de la hebra $+$ (azul) forman una montaña a la izquierda del sitio y las de la hebra $-$ (rojo), su imagen especular a la derecha (\cref{eq:13-hebras}); al desplazarlas $d/2$ hacia el centro (línea discontinua) se superponen sobre el sitio. (b) Correlación cruzada entre hebras (\cref{eq:13-xcor}) en un cromosoma simulado de \SI{3}{Mb} con regiones no mapeables. La muestra buena muestra un máximo claro en el tamaño de fragmento y un pico fantasma menor en la longitud de lectura ($L=36$); en la muestra pobre el pico fantasma domina, señal de que el anticuerpo apenas enriqueció. Datos: \archivo{figuras/cap13/generar.py}.}
\label{fig:13-hebras}
\end{figure}

\subsection{Llamar picos: el fondo de Poisson local}

Una vez desplazadas las lecturas, la pregunta es estadística: ¿en qué ventanas del genoma hay \emph{más} lecturas de las que produciría el fondo? El modelo más simple supone que, en ausencia de unión, las $N$ lecturas se reparten uniformemente sobre las $G$ posiciones mapeables.

\begin{teorema}{Conteo en una ventana bajo el fondo uniforme}{13-poisson}
Si cada una de las $N$ lecturas cae en una ventana de ancho $w$ con probabilidad $p=w/G$, de forma independiente, el número $X$ de lecturas en la ventana sigue una distribución binomial $\mathrm{Bin}(N,p)$. Cuando $N\to\infty$ y $p\to0$ con $Np=\lambda$ fijo,
\begin{equation}\label{eq:13-poisson}
\Prob(X=k) \;\longrightarrow\; \frac{e^{-\lambda}\lambda^{k}}{k!},\qquad \lambda=\frac{N\,w}{G},
\end{equation}
y el valor $p$ de observar al menos $k$ lecturas es la cola superior
\begin{equation}\label{eq:13-pvalor}
p(k) = \Prob(X\ge k) = 1-\sum_{i=0}^{k-1}\frac{e^{-\lambda}\lambda^{i}}{i!}.
\end{equation}
\end{teorema}

\begin{simbolos}
\simbolo{$N$}{Número total de lecturas (o fragmentos) mapeadas en la muestra tratada.}
\simbolo{$G$}{Tamaño efectivo del genoma: número de posiciones donde una lectura puede mapearse de forma única (unos $2{,}7\times10^9$ en humano).}
\simbolo{$w$}{Ancho de la ventana; en MACS, $w=2d$.}
\simbolo{$\lambda$}{Número esperado de lecturas en la ventana bajo el fondo.}
\simbolo{$k$}{Número de lecturas observadas en la ventana candidata.}
\end{simbolos}

La demostración es el límite clásico de Poisson:
\[
\binom{N}{k}p^k(1-p)^{N-k} = \frac{N(N-1)\cdots(N-k+1)}{N^k}\,\frac{\lambda^k}{k!}\,\Bigl(1-\frac{\lambda}{N}\Bigr)^{N-k},
\]
y el primer factor tiende a 1 mientras que el último tiende a $e^{-\lambda}$. Es la misma aproximación que, en el \cref{cap:06}, nos dio la teoría de Lander-Waterman para la cobertura.

El fondo uniforme, sin embargo, es falso. La sonicación rompe mejor la cromatina abierta; las regiones con más copias en la línea celular producen más lecturas; algunas zonas acumulan señal en cualquier experimento. Aplicar un $\lambda$ global en esas regiones llena la lista de picos de falsos positivos. La solución de MACS (\ingles{Model-based Analysis of ChIP-Seq}), propuesta por \textcite{c13-zhang2008}, es un \emph{Poisson dinámico}: para cada ventana candidata se toma como fondo el máximo de varias estimaciones, de modo que el modelo es conservador allí donde el control indica un sesgo.

\begin{definicion}{Lambda local de MACS}{13-lambdalocal}
Para una ventana candidata de ancho $2d$,
\begin{equation}\label{eq:13-lambdalocal}
\lambda_{\text{local}} = \max\bigl(\lambda_{\text{BG}},\ \lambda_{1k},\ \lambda_{5k},\ \lambda_{10k}\bigr),
\qquad
\lambda_{w'} = c_{w'}\cdot\frac{2d}{w'}\cdot\frac{N_T}{N_C},
\end{equation}
donde $\lambda_{\text{BG}}=N_T\,2d/G$ es el fondo global y $c_{w'}$ es el número de lecturas del \emph{control} en una ventana de ancho $w'\in\{1, 5, 10\}$ kb centrada en el candidato.
\end{definicion}

\begin{simbolos}
\simbolo{$\lambda_{\text{BG}}$}{Fondo global: las lecturas tratadas repartidas uniformemente.}
\simbolo{$c_{w'}$}{Lecturas del control en la ventana ancha $w'$; se reescalan al ancho $2d$ y a la profundidad del tratamiento.}
\simbolo{$N_T,\ N_C$}{Número de lecturas del tratamiento y del control.}
\end{simbolos}

Tomar el \emph{máximo} tiene una lógica de protección: si cualquiera de las escalas indica que la región es ruidosa, el umbral sube. Las versiones posteriores, MACS2 y MACS3, conservan esta idea (por omisión con ventanas de $d$, 1 kb y 10 kb en el control) y añaden la corrección de Benjamini-Hochberg para obtener valores $q$, el modo de picos anchos para marcas como H3K27me3 y el uso directo de fragmentos pareados. El artículo original estimaba además la tasa de falsos descubrimientos de forma empírica, intercambiando los papeles de tratamiento y control: los ``picos'' que aparecen en el control frente al tratamiento dan una idea de cuántos falsos positivos produce el procedimiento con el mismo umbral \parencite{c13-zhang2008}.

\begin{ejemplo}{Un pico, dos fondos}{13-macs}
Un experimento tiene $N_T=2\times10^7$ lecturas tratadas y $N_C=2{,}5\times10^7$ de control; el tamaño efectivo del genoma es $G=2{,}7\times10^9$ y $d=200$, así que la ventana mide 400 pb. El fondo global es
\[
\lambda_{\text{BG}} = \frac{2\times10^7\times400}{2{,}7\times10^9} = 2{,}963 .
\]
En el control hay 10 lecturas en la ventana de 1 kb, 45 en la de 5 kb y 95 en la de 10 kb. Reescalando con el factor de profundidad $N_T/N_C=0{,}8$: $\lambda_{1k}=10\cdot0{,}4\cdot0{,}8=3{,}200$, $\lambda_{5k}=45\cdot0{,}08\cdot0{,}8=2{,}880$ y $\lambda_{10k}=95\cdot0{,}04\cdot0{,}8=3{,}040$. Entonces $\lambda_{\text{local}}=3{,}200$. Con $k=17$ lecturas en la ventana,
\[
p = \Prob(X\ge17\mid\lambda=3{,}2) = 5{,}38\times10^{-8},\qquad -\log_{10}p = 7{,}27,
\]
un pico claro. Con el umbral de MACS por omisión en su versión original ($p<10^{-5}$), bastan 14 lecturas.

Ahora considere otra región, idéntica en el tratamiento, pero en una amplificación de la línea celular: el control tiene 400 lecturas en 10 kb, así que $\lambda_{10k}=12{,}8$. Con $k=14$ lecturas, el fondo global daría $p=2{,}96\times10^{-6}$ (un ``pico''), mientras que el fondo local da $p=0{,}405$: nada notable. El control ha evitado un falso positivo que ninguna profundidad de secuenciación habría corregido.
\end{ejemplo}

\begin{figure}[t]
\centering
\begin{tikzpicture}
\begin{groupplot}[group style={group size=2 by 1, horizontal sep=1.35cm},
   curso, height=5.3cm]
\nextgroupplot[width=7.4cm, title={(a) \ingles{pileup} y fondos},
   xlabel={posición (kb)}, ylabel={fragmentos por ventana de $2d$},
   xmin=0, xmax=10, ymin=0, ymax=20,
   legend style={at={(0.02,0.98)},anchor=north west}]
  \addplot[azul, thick] table[x=x,y=pile] {figuras/cap13/pileup.dat};
  \addlegendentry{tratamiento}
  \addplot[gris, thick] table[x=x,y=ctrl] {figuras/cap13/pileup.dat};
  \addlegendentry{control reescalado}
  \addplot[naranja, very thick] table[x=x,y=laml] {figuras/cap13/pileup.dat};
  \addlegendentry{$\lambda_{\text{local}}$}
  \addplot[rojo, dashed, thick, domain=0:10] {2.963};
  \addlegendentry{$\lambda_{\text{BG}}$}
  \node[etiqueta,anchor=west] at (axis cs:5.45,17.5) {pico};
  \node[etiqueta,anchor=south] at (axis cs:2.3,6.9) {débil};
\nextgroupplot[width=5.7cm, title={(b) $\mathrm{Poisson}(\lambda=3{,}2)$},
   xlabel={lecturas en la ventana $k$}, ylabel={$\Prob(X=k)$},
   ymode=log, log origin=infty, ymin=1e-10, ymax=1,
   xmin=-0.8, xmax=20.8, ybar, bar width=3.2pt,
   ytick={1e-10,1e-8,1e-6,1e-4,1e-2,1}]
  \addplot[fill=azul!60, draw=azul] table[x=k,y=pmf,restrict x to domain=0:16] {figuras/cap13/pmf.dat};
  \addplot[fill=rojo!70, draw=rojooscuro] table[x=k,y=pmf,restrict x to domain=17:20] {figuras/cap13/pmf.dat};
  \node[etiqueta,anchor=south,text=rojooscuro] at (axis cs:18.5,1e-6) {$k\ge17$};
\end{groupplot}
\end{tikzpicture}
\caption[Llamada de picos con Poisson local]{(a) Diez kilobases simuladas: el \ingles{pileup} del tratamiento (azul) muestra un pico fuerte cerca de 5,2 kb y un enriquecimiento débil cerca de 2,3 kb. La línea naranja es el fondo local estimado del control suavizado, que nunca baja del fondo global (roja discontinua), tal como impone el máximo de la \cref{eq:13-lambdalocal}. (b) Distribución de Poisson con $\lambda_{\text{local}}=3{,}2$ en escala logarítmica. Diecisiete lecturas en una ventana de 400 pb caen en la cola roja, cuya probabilidad total es $5{,}4\times10^{-8}$ (\cref{ej:13-macs}). Observe cuán rápido decae la cola: cada lectura adicional multiplica la evidencia.}
\label{fig:13-poisson}
\end{figure}

\begin{consola}[Llamada de picos con MACS3]
# ChIP de un factor de transcripción, pares de lecturas
macs3 callpeak -t chip_rep1.bam -c input.bam \
    -f BAMPE -g hs -n ctcf_rep1 -q 0.01 --outdir picos/
# salidas: *_peaks.narrowPeak (BED6+4) y *_summits.bed
\end{consola}


\subsection{Calidad y reproducibilidad}

Una lista de picos no vale nada si no se sabe cuánto de ella es señal. Tres controles se han vuelto estándar.

\paragraph{FRiP} La fracción de lecturas en picos (\ingles{fraction of reads in peaks}) es el cociente entre las lecturas que caen dentro de los picos llamados y el total. Mide la relación señal/ruido de todo el experimento: si 3,1 millones de 20 millones de lecturas caen en picos, $\mathrm{FRiP}=0{,}155$. Las guías de ENCODE consideran aceptables, para factores de transcripción, valores del \SI{1}{\percent} o más, aunque los buenos experimentos con factores abundantes como CTCF superan con holgura esa cifra \parencite{c13-landt2012}.

\paragraph{Lista negra} Algunas regiones del genoma (repeticiones satélite cerca de centrómeros y telómeros, ADN ribosómico, secuencias con copias no representadas en la referencia) producen señal alta en \emph{cualquier} experimento de secuenciación, independientemente del tipo celular o del anticuerpo. \textcite{c13-amemiya2019} definieron la \emph{lista negra de ENCODE} para humano, ratón, gusano y mosca, y recomiendan eliminar sus regiones como paso de calidad esencial antes de cualquier análisis.

\paragraph{IDR} Dos réplicas biológicas de un mismo experimento deberían coincidir en sus picos fuertes y discrepar en los débiles, que son ruido. \textcite{c13-li2011} formalizaron esta intuición en la \emph{tasa de irreproducibilidad} (\ingles{irreproducible discovery rate}, IDR). Su modelo supone que cada pico pertenece a uno de dos grupos: reproducible, cuyas puntuaciones en las dos réplicas están correlacionadas, o irreproducible, cuyas puntuaciones son independientes. Las puntuaciones se transforman en rangos (para no depender de la escala de cada llamador de picos) y una mezcla de dos cópulas gaussianas se ajusta por EM, el mismo algoritmo que veremos en la \cref{sec:13-motivos}. Para cada pico $i$ se obtiene la probabilidad posterior de ser irreproducible, $\mathrm{idr}_i$, y para una lista de los picos más fuertes $S_\gamma$,
\begin{equation}\label{eq:13-idr}
\mathrm{IDR}(\gamma) = \frac{1}{|S_\gamma|}\sum_{i\in S_\gamma}\mathrm{idr}_i ,
\end{equation}
\begin{simbolos}
\simbolo{$\mathrm{idr}_i$}{Probabilidad posterior, según la mezcla, de que el pico $i$ pertenezca al grupo irreproducible.}
\simbolo{$S_\gamma$}{Conjunto de picos seleccionados con el umbral $\gamma$ (p.\,ej., los que tienen $\mathrm{idr}_i\le\gamma$).}
\end{simbolos}
\noindent que es el análogo exacto de la tasa de falsos descubrimientos: la fracción esperada de picos irreproducibles en la lista. El protocolo de ENCODE llama picos de forma permisiva en cada réplica y conserva los que superan el umbral de IDR elegido (típicamente 0,05) \parencite{c13-landt2012}. En una simulación con \num{1500} picos reales y \num{2500} de ruido, la correlación de Spearman global entre réplicas es de solo 0,70, pero los 857 picos que están entre los mil primeros de \emph{ambas} réplicas son todos reales: la estructura que el modelo aprende está en la cabeza de las listas.


\begin{cuidado}[Errores frecuentes en ChIP-seq]
\begin{enumerate}
  \item \emph{Sin control o con un control poco profundo.} El control debe tener al menos tanta profundidad como el tratamiento; si no, $\lambda_{\text{local}}$ se estima con mucho ruido.
  \item \emph{Duplicados de PCR.} Muchas lecturas idénticas apiladas en la misma posición inflan los conteos y violan la independencia del modelo de Poisson; se marcan y se eliminan (o se limita su número) antes de llamar picos.
\end{enumerate}
\end{cuidado}

\subsection{ATAC-seq: medir la accesibilidad}

ChIP-seq necesita un anticuerpo por proteína y millones de células. Una pregunta más general (¿qué regiones del genoma están \emph{abiertas}, libres de nucleosomas y disponibles para cualquier factor?) se responde con una idea ingeniosa. La transposasa Tn5, en una versión hiperactiva cargada con adaptadores de secuenciación, corta el ADN e inserta los adaptadores en el mismo acto (``tagmentación''). En la cromatina nativa, Tn5 solo accede a los segmentos no protegidos: las regiones libres de nucleosomas y los conectores entre nucleosomas. \textcite{c13-buenrostro2013} llamaron a la técnica ATAC-seq (\ingles{assay for transposase-accessible chromatin using sequencing}) y mostraron que un protocolo de dos pasos con 500 a \num{50000} células revela a la vez la cromatina abierta, la unión de factores y la posición de nucleosomas individuales; el protocolo detallado está en \textcite{c13-buenrostro2015}. Como Tn5 corta como dímero con un escalonamiento de 9 pb, los análisis de alta resolución desplazan las lecturas $+4/-5$ pb para situar el centro exacto de cada inserción \parencite{c13-buenrostro2013}.

\begin{figure}[t]
\centering
\begin{tikzpicture}[font=\sffamily\scriptsize,x=1cm,y=1cm]
  \tikzset{nucl/.style={circle,fill=gris!45,draw=gris,minimum size=7.5mm,inner sep=0pt},
           tn/.style={regular polygon,regular polygon sides=3,fill=magenta,draw=magenta!60!black,minimum size=3.2mm,inner sep=0pt,shape border rotate=180}}
  \draw[azul,thick] (0,1.5) -- (12.8,1.5);
  \foreach \x in {1.0,2.25,3.5,8.3,9.55,10.8,12.05}{\node[nucl] at (\x,1.5) {};}
  \node[etiqueta] at (5.9,2.35) {región libre de nucleosomas};
  \draw[naranja,thick] (4.6,2.1) -- (7.2,2.1);
  \node[fill=naranja,draw=naranja!70!black,rounded corners=1.5pt,minimum width=5mm,minimum height=3mm] at (5.9,1.72) {};
  \node[etiqueta,text=naranja!70!black] at (5.9,1.05) {factor};
  % inserciones Tn5
  \foreach \x in {1.625,2.875,4.4,5.25,6.55,7.4,8.925,10.175,11.425}{
     \node[tn] at (\x,2.05) {};
     \draw[magenta,densely dotted] (\x,1.9) -- (\x,1.55);}
  \node[etiqueta,text=magenta!70!black,anchor=west] at (12.2,2.25) {Tn5};
  % fragmentos resultantes
  \draw[{Bar[width=4pt]}-{Bar[width=4pt]},aqua!70!black,thick] (4.4,0.45) -- node[below,etiqueta,text=aqua!40!black]{$<$100 pb (NFR)} (5.25,0.45);
  \draw[{Bar[width=4pt]}-{Bar[width=4pt]},aqua!70!black,thick] (6.55,0.45) -- (7.4,0.45);
  \draw[{Bar[width=4pt]}-{Bar[width=4pt]},azul,thick] (8.925,0.45) -- node[below,etiqueta,text=azulprofundo]{$\sim$200 pb (mono)} (10.175,0.45);
  \draw[{Bar[width=4pt]}-{Bar[width=4pt]},violeta,thick] (1.625,-0.15) -- node[below,etiqueta,text=violeta]{$\sim$400 pb (di)} (4.4,-0.15);
\end{tikzpicture}
\caption[Principio de ATAC-seq]{Principio de ATAC-seq. La transposasa Tn5 (triángulos magenta) inserta adaptadores solo donde el ADN no está protegido: en la región libre de nucleosomas (NFR) alrededor de un factor unido y en los conectores entre nucleosomas (círculos grises). Dos inserciones sucesivas definen un fragmento secuenciable; su longitud revela lo que había entre ellas: nada (fragmentos cortos de la NFR), un nucleosoma (unos 200 pb) o dos (unos 400 pb).}
\label{fig:13-tn5}
\end{figure}

Dos inserciones de Tn5 definen un fragmento, y su longitud es un registro directo de lo que había entre ellas (\cref{fig:13-tn5}). Por eso la distribución de tamaños de fragmento de un buen experimento de ATAC-seq tiene una forma inconfundible: una gran población por debajo de 100 pb, procedente de regiones libres de nucleosomas, y una serie de ``jorobas'' en múltiplos de la repetición nucleosomal. Podemos escribirla como una mezcla:
\begin{equation}\label{eq:13-atac}
h(\ell) \;=\; \pi_0\,h_0(\ell) \;+\; \sum_{j=1}^{J}\pi_j\,\mathcal N\!\left(\ell;\ j\,\ell_{\text{nuc}},\ \sigma_j^2\right),
\qquad \sum_{j=0}^{J}\pi_j = 1 ,
\end{equation}
\begin{simbolos}
\simbolo{$h(\ell)$}{Densidad de la longitud $\ell$ de los fragmentos.}
\simbolo{$h_0$}{Componente libre de nucleosomas (fragmentos cortos, con moda por debajo de 100 pb).}
\simbolo{$\ell_{\text{nuc}}$}{Repetición nucleosomal efectiva (unos 190--200 pb: 147 del núcleo más el conector).}
\simbolo{$\pi_j,\ \sigma_j$}{Proporción y dispersión de los fragmentos que abarcan $j$ nucleosomas.}
\end{simbolos}
\noindent A esta forma gruesa se superpone una oscilación fina de unos 10,5 pb, el paso de la doble hélice: Tn5 corta preferentemente el lado de la hélice expuesto, de modo que las distancias entre cortes favorecen múltiplos de una vuelta. La \cref{fig:13-atac}\,(a) muestra una distribución simulada con ese modelo, en la que el \SI{38}{\percent} de los fragmentos mide menos de 100 pb y el \SI{21}{\percent} cae en el intervalo mononucleosomal (180--247 pb).

\begin{definicion}{Enriquecimiento en TSS}{13-tss}
Sea $I(x)$ el número de inserciones de Tn5 acumuladas en la posición $x$ relativa a todos los sitios de inicio de la transcripción (TSS) del genoma, orientados según la hebra del gen. El perfil normalizado es
\begin{equation}\label{eq:13-tss}
E(x) = \frac{I(x)}{\tfrac{1}{2m}\left(\sum_{u=-2000}^{-2000+m-1} I(u) + \sum_{u=2000-m+1}^{2000} I(u)\right)},
\end{equation}
es decir, el conteo dividido por el promedio en los extremos de la ventana ($m$ posiciones a cada lado, típicamente 100 pb). La \emph{puntuación de enriquecimiento en TSS} es $\max_x E(x)$ cerca de $x=0$.
\end{definicion}

\begin{simbolos}
\simbolo{$I(x)$}{Inserciones acumuladas a distancia $x$ del TSS (en pb, con signo según la orientación del gen).}
\simbolo{$m$}{Ancho de los bordes usados como fondo.}
\simbolo{$E(x)$}{Enriquecimiento relativo al fondo lejano; vale $\approx1$ lejos del TSS.}
\end{simbolos}

Los promotores activos son las regiones abiertas más universales, así que un experimento que ``ve'' cromatina abierta debe mostrar un pico agudo en el TSS. Un enriquecimiento bajo (digamos, inferior a 5) indica demasiadas lecturas de fondo, típicamente por células muertas o lisis excesiva de los núcleos. El perfil simulado de la \cref{fig:13-atac}\,(b) alcanza 8,35, con un pequeño hombro río abajo que corresponde al nucleosoma $+1$, bien posicionado en los genes activos.

\begin{figure}[t]
\centering
\begin{tikzpicture}
\begin{groupplot}[group style={group size=2 by 1, horizontal sep=1.35cm},
   curso, height=5.4cm]
\nextgroupplot[width=7.3cm, title={(a) tamaños de fragmento},
   xlabel={longitud del fragmento (pb)}, ylabel={frecuencia (\textperthousand{} por pb)},
   ymode=log, xmin=0, xmax=1000, ymin=0.01, ymax=15,
   ytick={0.01,0.1,1,10}, yticklabels={0{,}01,0{,}1,1,10}]
  \addplot[azul, thick] table[x=l,y=f] {figuras/cap13/atac_frag.dat};
  \node[etiqueta,text=aqua!40!black] at (axis cs:100,0.6) {NFR};
  \node[etiqueta,anchor=south] at (axis cs:205,6) {mono};
  \node[etiqueta,anchor=south] at (axis cs:390,1.9) {di};
  \node[etiqueta,anchor=south] at (axis cs:580,0.55) {tri};
\nextgroupplot[width=5.8cm, title={(b) perfil en el TSS},
   xlabel={distancia al TSS (pb)}, ylabel={enriquecimiento $E(x)$},
   xmin=-2000, xmax=2000, ymin=0, ymax=9.5, xtick={-2000,-1000,0,1000,2000},
   xticklabels={$-2$k,$-1$k,0,1k,2k}]
  \addplot[naranja, very thick] table[x=x,y=e] {figuras/cap13/tss.dat};
  \draw[tinta2,dashed] (axis cs:-2000,1) -- (axis cs:2000,1);
  \node[etiqueta,anchor=west] at (axis cs:180,8.3) {máx.\ 8,35};
  \node[etiqueta,anchor=west,align=left] at (axis cs:420,3.2) {nucleosoma\\$+1$};
  \draw[tinta2,-{Stealth}] (axis cs:420,2.9) -- (axis cs:340,1.8);
\end{groupplot}
\end{tikzpicture}
\caption[Controles de calidad de ATAC-seq]{Controles de calidad de ATAC-seq sobre datos simulados con la mezcla de la \cref{eq:13-atac}. (a) Distribución de tamaños de fragmento en escala logarítmica: la población libre de nucleosomas (NFR) domina, y le siguen jorobas mono-, di- y trinucleosomales separadas por unos 190 pb; la ondulación fina de 10,5 pb refleja la periodicidad helicoidal del ADN. Una biblioteca sin jorobas sugiere digestión excesiva o ADN desnudo. (b) Enriquecimiento de inserciones alrededor del TSS (\cref{eq:13-tss}), normalizado por los extremos de la ventana (línea discontinua en 1).}
\label{fig:13-atac}
\end{figure}


\paragraph{CUT\&RUN} Una alternativa moderna a ChIP-seq invierte la lógica del experimento. En lugar de fragmentar toda la cromatina y luego pescar los trozos que interesan, CUT\&RUN (\ingles{cleavage under targets and release using nuclease}) lleva una nucleasa hasta la proteína de interés: un anticuerpo se une al factor en núcleos intactos y una nucleasa micrococal acoplada a proteína A corta el ADN a su alrededor, liberando solo los fragmentos unidos. Como no hay entrecruzamiento ni fragmentación global, el fondo es mínimo: \textcite{c13-skene2017} obtuvieron perfiles precisos con aproximadamente una décima parte de la profundidad de secuenciación de ChIP-seq.

\begin{ideaclave}
Un pico es un exceso de lecturas sobre un fondo \emph{local}. La geometría de hebras delata la señal verdadera, el control protege frente a sesgos del genoma y la reproducibilidad entre réplicas, no el valor $p$, decide qué picos se conservan.
\end{ideaclave}

% ---------------------------------------------------------------------
\section{Metilación del ADN}\label{sec:13-metil}
% ---------------------------------------------------------------------
\enlacecolab{modulo-13-epigenomica/13.2_metilacion_adn.ipynb}{13.2 · Metilación del ADN}

\begin{intuicion}[Post-its en un libro de recetas]
Piense en un libro de recetas compartido por toda una familia. Nadie reescribe las recetas, pero cada cocinero pega notas adhesivas sobre las que no quiere que se usen: ``no preparar'', ``solo en Navidad''. Las notas no cambian el texto, se copian cuando alguien fotocopia el libro para un hijo, y pueden despegarse. La metilación del ADN funciona de manera parecida: un grupo metilo sobre una citosina no altera la secuencia, se hereda en cada división celular gracias a una enzima que copia el patrón de la hebra vieja a la nueva, y puede borrarse. El problema técnico es que las fotocopias normales, es decir, la PCR y la secuenciación, no copian las notas. Hay que encontrar una tinta que revele dónde estaban antes de fotocopiar.
\end{intuicion}

\subsection{Citosinas metiladas e islas CpG}

En los mamíferos, la metilación del ADN consiste casi exclusivamente en añadir un grupo metilo al carbono 5 de una citosina (5mC) seguida de una guanina, el dinucleótido CpG. Como el dinucleótido complementario de CpG es también CpG, la marca es simétrica en las dos hebras, y eso permite heredarla: tras la replicación, la hebra nueva se metila frente a los CpG metilados de la hebra vieja. La metilación de un promotor se asocia en general al silenciamiento estable del gen, y participa en la impronta genómica, la inactivación del cromosoma X y el control de elementos transponibles.

La metilación deja una huella en la propia secuencia del genoma. La 5mC tiende a desaminarse espontáneamente a timina, un error que la reparación corrige peor que la desaminación de la citosina normal a uracilo. A lo largo de la evolución, los CpG metilados se han ido convirtiendo en TpG (y CpA en la otra hebra), de modo que el genoma humano tiene alrededor de una cuarta parte de los CpG que esperaríamos por su composición. La excepción son las \emph{islas CpG}: tramos de cientos a miles de pares de bases con abundantes CpG no metilados, que \textcite{c13-bird1986} asoció con los extremos 5' de la mayoría de los genes de vertebrados y, en particular, de los genes de mantenimiento, cuya disponibilidad debe ser constante en el núcleo. \textcite{c13-gardinergarden1987} dieron el criterio operativo que todavía se usa, basado en la razón entre CpG observados y esperados.

\begin{definicion}{Isla CpG (criterio de Gardiner-Garden y Frommer)}{13-islacpg}
Una región de longitud $L\ge200$ pb es una isla CpG si su contenido de GC supera el \SI{50}{\percent} y
\begin{equation}\label{eq:13-oe}
\frac{\text{CpG obs.}}{\text{CpG esp.}} = \frac{n_{CG}}{n_C\,n_G/L} = \frac{n_{CG}\,L}{n_C\,n_G} > 0{,}6 .
\end{equation}
\end{definicion}

\begin{simbolos}
\simbolo{$n_{CG}$}{Número de dinucleótidos CpG en la región (en una hebra).}
\simbolo{$n_C,\ n_G$}{Número de citosinas y guaninas.}
\simbolo{$L$}{Longitud de la región; $n_Cn_G/L$ es el número de CpG esperado si C y G se sucedieran al azar.}
\end{simbolos}

\begin{ejemplo}{Dos regiones de 100 pb}{13-oe}
Generamos dos secuencias sintéticas de 100 pb que imitan la evolución de la desaminación: una rica en GC en la que se conserva el \SI{90}{\percent} de los CpG (``isla'') y otra con GC moderado en la que solo se conserva el \SI{22}{\percent} (``mar''). En la primera hay $n_C=37$, $n_G=32$ y $n_{CG}=12$: GC $=\SI{69}{\percent}$ y
\[
\frac{\text{obs.}}{\text{esp.}} = \frac{12\times100}{37\times32} = 1{,}01 > 0{,}6 .
\]
En la segunda, $n_C=22$, $n_G=17$ y un único CpG: GC $=\SI{39}{\percent}$ y obs./esp. $=100/(22\times17)=0{,}27$, el valor típico del ``mar'' genómico. (Para aplicar formalmente el criterio habría que usar ventanas de al menos 200 pb.)
\end{ejemplo}

\subsection{La química del bisulfito}

Ni la PCR ni los secuenciadores de síntesis distinguen 5mC de C: la polimerasa pone una G frente a ambas y el grupo metilo se pierde en la primera copia. La solución clásica es química. El bisulfito de sodio desamina la citosina a uracilo en condiciones en las que la 5-metilcitosina, protegida por el metilo, apenas reacciona. \textcite{c13-frommer1992} convirtieron esta reacción en un método de secuenciación genómica que identifica positivamente cada 5mC: tras el tratamiento, se amplifica por PCR y en el producto todos los uracilos (y timinas) aparecen como T, mientras que solo las 5mC siguen leyéndose como C (\cref{fig:13-bisulfito}).

\begin{figure}[t]
\centering
\begin{tikzpicture}[font=\sffamily\scriptsize]
  \tikzset{bU/.style={minimum size=5.5mm, inner sep=0pt, font=\ttfamily\bfseries\small, text=white, rounded corners=1.5pt, fill=magenta},
           bm/.style={nt=C, draw=azulnoche, line width=1.2pt}}
  \def\xs{0.6}\def\capbM{M}\def\capbU{U}
  % fila 1: original
  \node[anchor=east,etiqueta,align=right] at (-0.25,3.0) {ADN genómico};
  \foreach \b [count=\i from 0] in {A,M,G,T,C,G,A,C,T,T,M,G,A,C}{
     \ifx\b\capbM \node[bm] at (\i*\xs,3.0) {C}; \fill[naranja] (\i*\xs+0.2,3.27) circle (0.07);
     \else \node[nt=\b] at (\i*\xs,3.0) {\b}; \fi}
  % fila 2: bisulfito
  \node[anchor=east,etiqueta,align=right] at (-0.25,2.0) {tras bisulfito};
  \foreach \b [count=\i from 0] in {A,M,G,T,U,G,A,U,T,T,M,G,A,U}{
     \ifx\b\capbM \node[bm] at (\i*\xs,2.0) {C};
     \else\ifx\b\capbU \node[bU] at (\i*\xs,2.0) {U};
     \else \node[nt=\b] at (\i*\xs,2.0) {\b}; \fi\fi}
  % fila 3: PCR
  \node[anchor=east,etiqueta,align=right] at (-0.25,1.0) {tras PCR (lectura)};
  \foreach \b [count=\i from 0] in {A,C,G,T,T,G,A,T,T,T,C,G,A,T}{
     \node[nt=\b] at (\i*\xs,1.0) {\b};}
  % fila 4: interpretación
  \node[anchor=east,etiqueta,align=right] at (-0.25,0.2) {lectura de la C};
  \foreach \i/\t/\c in {1/{met.}/azulprofundo,4/{no met.}/magenta!70!black,7/{no met.}/magenta!70!black,10/{met.}/azulprofundo,13/{no met.}/magenta!70!black}{
     \node[etiqueta,text=\c,font=\sffamily\tiny\bfseries] at (\i*\xs,0.2) {\t};
     \draw[\c,thin] (\i*\xs,0.7) -- (\i*\xs,0.38);}
  \foreach \y in {2.5,1.5}{\draw[flecha=tinta2] (8.4,\y+0.2) -- (8.4,\y-0.2);}
  \node[etiqueta,anchor=west,align=left] at (8.55,2.5) {C $\to$ U\\(5mC resiste)};
  \node[etiqueta,anchor=west,align=left] at (8.55,1.5) {U $\to$ T\\(5mC $\to$ C)};
  \node[etiqueta,anchor=west,align=left,text=azulprofundo] at (8.55,3.0) {\textcolor{naranja}{$\bullet$} 5mC};
\end{tikzpicture}
\caption[Conversión por bisulfito]{Química de la conversión por bisulfito en una hebra de 14 bases con dos CpG metilados (C con borde oscuro y punto naranja), un CpG no metilado (posición 5) y dos citosinas fuera de contexto CpG. El bisulfito desamina las citosinas no metiladas a uracilo (magenta); la PCR copia el uracilo como timina. Al comparar la lectura con la referencia, una C que sigue siendo C estaba metilada; una C que aparece como T no lo estaba. La hebra convertida ya no es complementaria de la original, y por eso cada hebra del genoma produce lecturas distintas.}
\label{fig:13-bisulfito}
\end{figure}

La consecuencia para el análisis es doble. Primero, la medida es \emph{digital} en cada molécula: cada lectura dice ``C'' o ``T'' en cada citosina de la referencia, y el nivel de metilación de un sitio es la fracción de lecturas que dicen ``C''. Segundo, el alineamiento se complica, porque una lectura convertida difiere de la referencia en casi todas sus citosinas. Si la conversión fuera incompleta, algunas C no metiladas se leerían como metiladas; por eso se estima la tasa de no conversión con ADN control sin metilar añadido a la muestra, o con citosinas fuera de contexto CpG en tejidos donde apenas se metilan.

\subsection{Secuenciación con bisulfito del genoma completo y alineamiento}

La secuenciación con bisulfito del genoma completo (WGBS, \ingles{whole-genome bisulfite sequencing}) aplica la conversión a una biblioteca de todo el genoma. \textcite{c13-lister2009} publicaron así los primeros metilomas humanos con resolución de una base, de células madre embrionarias y fibroblastos fetales, y descubrieron que casi una cuarta parte de la metilación de las células madre embrionarias estaba fuera del contexto CG, una forma que desaparecía al diferenciarse.

\paragraph{El problema del alineamiento} Un alineador convencional penalizaría cada T de la lectura frente a una C de la referencia. Pero no basta con permitir esos desajustes sin coste: tras la conversión, la lectura está escrita casi en un alfabeto de tres letras, y su complejidad baja tanto que muchas lecturas mapearían en varios lugares. La solución de Bismark \parencite{c13-krueger2011} es convertir \emph{todo} de forma computacional antes de alinear (\cref{fig:13-bismark}): en las lecturas se reemplazan todas las C por T, y el genoma se prepara en dos versiones, una con todas las C convertidas en T (para la hebra superior) y otra con todas las G convertidas en A (para la inferior). Cada lectura se alinea con un alineador estándar contra las versiones convertidas en todas las combinaciones posibles de hebra; se conserva solo el alineamiento único con menos desajustes y, una vez ubicado, se compara la lectura \emph{original} con la referencia \emph{original} para llamar la metilación de cada citosina, distinguiendo los contextos CpG, CHG y CHH (con H = A, C o T).

\begin{figure}[t]
\centering
\begin{tikzpicture}[font=\sffamily\scriptsize, node distance=4mm]
  \node[caja=azul, text width=2.1cm, font=\sffamily\scriptsize] (lec) at (0,1.4) {lectura original\\\texttt{ATTGATTTCGAT}};
  \node[caja=azul, text width=2.1cm, font=\sffamily\scriptsize] (lct) at (0,0) {lectura C$\to$T\\\texttt{ATTGATTTTGAT}};
  \node[caja=naranja, text width=2.2cm, font=\sffamily\scriptsize] (g1) at (3.55,1.4) {genoma C$\to$T\\(hebra superior)};
  \node[caja=naranja, text width=2.2cm, font=\sffamily\scriptsize] (g2) at (3.55,0) {genoma G$\to$A\\(hebra inferior)};
  \node[caja=violeta, text width=2.2cm, font=\sffamily\scriptsize] (al) at (6.95,0.7) {alinear en todas las\\combinaciones;\\quedarse con el\\mejor único};
  \node[caja=aqua, text width=2.3cm, font=\sffamily\scriptsize] (ll) at (10.3,0.7) {comparar lectura y\\referencia \emph{originales}:\\CpG, CHG, CHH};
  \draw[flecha] (lec) -- (lct);
  \draw[flecha] (lct.east) -- (g1.west);
  \draw[flecha] (lct.east) -- (g2.west);
  \draw[flecha] (g1.east) -- (al.west);
  \draw[flecha] (g2.east) -- (al.west);
  \draw[flecha] (al) -- (ll);
  \draw[flecha=azul, dashed, rounded corners=4pt] (lec.north) -- ++(0,0.35) -| (ll.north);
\end{tikzpicture}
\caption[Estrategia de alineamiento de Bismark]{Estrategia de alineamiento de Bismark \parencite{c13-krueger2011}. La lectura y el genoma se convierten \emph{in silico} al alfabeto reducido, de modo que las diferencias debidas al bisulfito desaparecen y solo quedan los verdaderos desajustes. Tras encontrar la posición única, la metilación se llama volviendo a la lectura original (flecha discontinua): la C del CpG de la lectura de ejemplo indica una citosina metilada; las T frente a C de la referencia, citosinas no metiladas.}
\label{fig:13-bismark}
\end{figure}

\begin{consola}[Flujo básico de WGBS con Bismark]
bismark_genome_preparation genoma/            # versiones C->T y G->A
bismark --genome genoma/ -1 m1_R1.fq.gz -2 m1_R2.fq.gz
deduplicate_bismark --bam m1_R1_bismark_bt2_pe.bam
bismark_methylation_extractor --bedGraph --cytosine_report \
    --genome_folder genoma/ m1_R1_bismark_bt2_pe.deduplicated.bam
# salida: por cada CpG, lecturas metiladas y no metiladas
\end{consola}

\subsection{Cuantificar: valores $\beta$ y $M$}

Para un CpG cubierto por $n$ lecturas, de las cuales $c$ dicen ``C'', el modelo natural es binomial: $c\sim\mathrm{Bin}(n,\beta)$, donde $\beta$ es la fracción de moléculas metiladas en la muestra. El estimador de máxima verosimilitud es $\hat\beta=c/n$, con error estándar $\sqrt{\hat\beta(1-\hat\beta)/n}$.

Las micromatrices de Illumina (Infinium 450K y EPIC), aún muy usadas en estudios poblacionales, miden lo mismo de otra forma: para cada CpG hay dos intensidades de fluorescencia, $M_s$ para el alelo metilado y $U_s$ para el no metilado (tras la conversión, la sonda distingue C de T). \textcite{c13-du2010} compararon las dos métricas en uso.

\begin{definicion}{Valores $\beta$ y $M$}{13-betam}
\begin{equation}\label{eq:13-beta}
\beta = \frac{\max(M_s,0)}{\max(M_s,0)+\max(U_s,0)+\alpha},\qquad
M = \log_2\frac{\max(M_s,0)+\alpha'}{\max(U_s,0)+\alpha'} .
\end{equation}
Ignorando las constantes, $M=\log_2\dfrac{\beta}{1-\beta}$: el valor $M$ es la transformación \emph{logit} (en base 2) del valor $\beta$.
\end{definicion}

\begin{simbolos}
\simbolo{$M_s,\ U_s$}{Intensidades de las sondas metilada y no metilada del CpG $s$ (en WGBS, los conteos $c$ y $n-c$).}
\simbolo{$\alpha$}{Constante que estabiliza $\beta$ cuando ambas intensidades son bajas (Illumina usa 100).}
\simbolo{$\alpha'$}{Pequeña constante que evita el logaritmo de cero (p.\,ej., 1).}
\simbolo{$\beta\in[0,1]$}{Fracción de metilación; interpretable como porcentaje de moléculas metiladas.}
\simbolo{$M\in\R$}{Valor $M$; $M=0$ corresponde a $\beta=0{,}5$.}
\end{simbolos}

¿Por qué dos métricas? El valor $\beta$ tiene una interpretación biológica directa, pero su varianza depende de su media: cerca de 0 y de 1 está comprimido por los límites del intervalo. El método delta muestra que la transformación logit compensa exactamente esa compresión. Si $\operatorname{Var}(\hat\beta)\approx\beta(1-\beta)/n$ (binomial), entonces
\begin{equation}\label{eq:13-delta}
\operatorname{Var}(\hat M) \approx \left(\frac{dM}{d\beta}\right)^{2}\operatorname{Var}(\hat\beta)
= \frac{1}{\bigl(\beta(1-\beta)\ln2\bigr)^2}\cdot\frac{\beta(1-\beta)}{n}
= \frac{1}{n\,\beta(1-\beta)\,(\ln 2)^2},
\end{equation}
\begin{simbolos}
\simbolo{$dM/d\beta$}{Derivada de la transformación logit en base 2, $1/(\beta(1-\beta)\ln2)$.}
\simbolo{$n$}{Número de moléculas (lecturas) que sustentan la medida.}
\end{simbolos}
\noindent y los contrastes en la escala $M$ ya no tratan igual una diferencia de 0,05 en torno a $\beta=0{,}5$ que en torno a $\beta=0{,}97$, donde la segunda es mucho más significativa biológica y estadísticamente. \textcite{c13-du2010} concluyeron que $\beta$ es preferible para informar resultados y $M$ para el análisis diferencial en micromatrices, cuya varianza técnica es más homogénea en esa escala. En WGBS, donde los datos son conteos, lo más limpio es modelar directamente la binomial, como veremos enseguida.

\begin{figure}[t]
\centering
\begin{tikzpicture}
\begin{groupplot}[group style={group size=3 by 1, horizontal sep=1.25cm},
   curso, height=4.8cm]
\nextgroupplot[width=4.9cm, title={(a) valores $\beta$}, xlabel={$\beta$}, ylabel={densidad},
   xmin=0, xmax=1, ymin=0, ymax=5.2, xtick={0,0.5,1}]
  \addplot[const plot, fill=azul!40, draw=azul, thick] table[x=x,y=d] {figuras/cap13/beta_hist.dat} \closedcycle;
\nextgroupplot[width=4.9cm, title={(b) valores $M$}, xlabel={$M$}, ylabel={densidad},
   xmin=-8, xmax=6, ymin=0, ymax=0.25, xtick={-8,-4,0,4}]
  \addplot[const plot, fill=naranja!40, draw=naranja, thick] table[x=x,y=d] {figuras/cap13/m_hist.dat} \closedcycle;
\nextgroupplot[width=4.4cm, title={(c) $M=\log_2\frac{\beta}{1-\beta}$}, xlabel={$\beta$}, ylabel={$M$},
   xmin=0, xmax=1, ymin=-7, ymax=7, xtick={0,0.5,1}]
  \addplot[violeta, very thick, domain=0.008:0.992, samples=120] {ln(x/(1-x))/ln(2)};
  \draw[tinta2,dashed] (axis cs:0,0) -- (axis cs:1,0);
\end{groupplot}
\end{tikzpicture}
\caption[Distribución de valores $\beta$ y $M$]{Distribución de \num{400000} CpG simulados con intensidades de tipo Infinium. (a) La distribución de $\beta$ es fuertemente bimodal: el \SI{32,5}{\percent} de los sitios tiene $\beta<0{,}2$ (sobre todo CpG de islas, no metilados) y el \SI{43,1}{\percent} tiene $\beta>0{,}8$ (el mar genómico, metilado). (b) Los mismos datos en escala $M$: los dos modos se separan y se ensanchan. (c) La transformación logit estira los extremos del intervalo $[0,1]$, compensando la compresión de la varianza de $\beta$ cerca de 0 y de 1 (\cref{eq:13-delta}).}
\label{fig:13-beta}
\end{figure}

\subsection{Metilación diferencial: el modelo beta-binomial}

Al comparar dos grupos (tumor y tejido normal, tratados y controles) buscamos CpG individuales (\ingles{differentially methylated loci}, DML) o regiones (DMR) cuyo nivel de metilación difiera. El modelo binomial captura el muestreo de lecturas, pero no la variación \emph{biológica} entre réplicas: dos individuos sanos no tienen exactamente la misma $\beta$ en un CpG. Ignorar esa variabilidad subestima la varianza y produce avalanchas de falsos positivos, exactamente el mismo fenómeno de sobredispersión que en RNA-seq obliga a pasar de la Poisson a la binomial negativa (\cref{cap:11}).

\begin{teorema}{Varianza beta-binomial}{13-betabin}
Si en la réplica $r$ la fracción de metilación es $p_r\sim\mathrm{Beta}(a,b)$ con media $\mu=a/(a+b)$ y dispersión $\phi=1/(a+b+1)$, y $c_r\mid p_r\sim\mathrm{Bin}(n_r,p_r)$, entonces
\begin{equation}\label{eq:13-betabin}
\E[c_r]=n_r\,\mu,\qquad \operatorname{Var}(c_r) = n_r\,\mu(1-\mu)\,\bigl[1+(n_r-1)\,\phi\bigr].
\end{equation}
\end{teorema}

\begin{simbolos}
\simbolo{$c_r,\ n_r$}{Lecturas metiladas y cobertura total del CpG en la réplica $r$.}
\simbolo{$p_r$}{Fracción de metilación verdadera de la réplica (varía entre individuos).}
\simbolo{$\mu$}{Media de metilación del grupo en ese CpG.}
\simbolo{$\phi\in[0,1)$}{Dispersión: correlación entre dos lecturas de la misma réplica; $\phi=0$ recupera la binomial.}
\end{simbolos}

La demostración usa la ley de la varianza total: $\operatorname{Var}(c)=\E[\operatorname{Var}(c\mid p)]+\operatorname{Var}(\E[c\mid p]) = n\E[p(1-p)] + n^2\operatorname{Var}(p)$. Con $\E[p(1-p)] = \mu(1-\mu)-\operatorname{Var}(p)$ y $\operatorname{Var}(p)=\mu(1-\mu)\phi$ para la beta, se obtiene $n\mu(1-\mu)(1-\phi)+n^2\mu(1-\mu)\phi$, que es la \cref{eq:13-betabin}. El factor $1+(n-1)\phi$ tiene una lectura importante: \emph{aumentar la cobertura no elimina la variación biológica}. Con $n=20$ y $\mu=0{,}85$, la desviación estándar de $c$ es 1,60 si $\phi=0$, pero 2,23 con $\phi=0{,}05$ y 2,72 con $\phi=0{,}1$.

El paquete DSS de \textcite{c13-feng2014} usa este modelo con una idea adicional: como con pocas réplicas la dispersión de cada CpG se estima muy mal, le asigna una distribución a priori lognormal común a todos los CpG, de modo que cada estimación se ``encoge'' hacia el valor típico del genoma (compartir información entre sitios). Con la dispersión estimada, contrasta la diferencia de medias con una prueba de Wald:
\begin{equation}\label{eq:13-wald}
z = \frac{\hat\mu_A-\hat\mu_B}{\sqrt{\widehat{\operatorname{Var}}(\hat\mu_A)+\widehat{\operatorname{Var}}(\hat\mu_B)}},\qquad
\widehat{\operatorname{Var}}(\hat\mu_g)\approx \frac{1}{R_g^2}\sum_{r=1}^{R_g}\frac{\hat\mu_g(1-\hat\mu_g)\bigl[1+(n_r-1)\hat\phi\bigr]}{n_r},
\end{equation}
\begin{simbolos}
\simbolo{$\hat\mu_g$}{Metilación estimada del grupo $g$: $\sum_r c_r/\sum_r n_r$.}
\simbolo{$R_g$}{Número de réplicas del grupo $g$.}
\simbolo{$\hat\phi$}{Dispersión estimada (encogida) del CpG.}
\simbolo{$z$}{Estadístico de Wald; bajo $H_0$ sigue aproximadamente una normal estándar.}
\end{simbolos}

\begin{ejemplo}{Un CpG, tres réplicas por grupo}{13-wald}
Primero, una sola muestra: 17 de 20 lecturas dicen ``C''. Entonces $\hat\beta=0{,}850$, con error estándar binomial $0{,}080$. El intervalo de Wald al \SI{95}{\percent}, $(0{,}694;\ 1{,}006)$, se sale de $[0,1]$: con proporciones cercanas a los extremos conviene un intervalo bayesiano de Jeffreys (a priori $\mathrm{Beta}(\tfrac12,\tfrac12)$), que da $(0{,}651;\ 0{,}956)$. El valor $M$, con $\alpha'=1$, es $\log_2(18/4)=2{,}17$.

Ahora, el grupo A tiene conteos $17/20$, $22/25$ y $14/18$, y el grupo B, $6/21$, $9/24$ y $5/19$. Las medias son $\hat\mu_A=53/63=0{,}841$ y $\hat\mu_B=20/64=0{,}312$. Con $\hat\phi=0{,}05$, la \cref{eq:13-wald} da $\widehat{\operatorname{Var}}(\hat\mu_A)=0{,}00428$, $\widehat{\operatorname{Var}}(\hat\mu_B)=0{,}00680$ y
\[
z=\frac{0{,}841-0{,}312}{\sqrt{0{,}00428+0{,}00680}} = 5{,}02,\qquad p = 5{,}1\times10^{-7}.
\]
Si ignoráramos la dispersión ($\phi=0$) obtendríamos $z=7{,}10$ y $p=1{,}3\times10^{-12}$: cinco órdenes de magnitud de ``evidencia'' que no existen. Con $\phi=0{,}1$, $z=4{,}10$; con $\phi=0{,}2$, $z=3{,}18$ y $p=0{,}0015$. La conclusión depende críticamente de cuánto varían los individuos, y por eso las réplicas biológicas son irrenunciables.
\end{ejemplo}


\paragraph{De sitios a regiones} Los CpG vecinos están fuertemente correlacionados: la metilación cambia por bloques de cientos de pares de bases, no sitio a sitio. Los métodos de DMR aprovechan esa estructura suavizando los niveles de metilación a lo largo del genoma (lo que además ``presta'' cobertura de los vecinos a los CpG poco cubiertos) y agrupando CpG significativos contiguos en regiones, con requisitos mínimos de longitud, número de CpG y diferencia media. La \cref{fig:13-dmr} muestra una simulación de 55 CpG en \SI{3}{kb}: la prueba de Wald con $\phi=0{,}05$ declara significativos ($p<0{,}01$) 15 CpG, de los que 14 están dentro de la DMR verdadera (\SIrange{1100}{1900}{pb}); el tramo contiguo más largo de CpG significativos, de \num{1159} a \num{1835} pb, recupera la región casi exactamente.

\def\capdmrini{1159}\def\capdmrfin{1835}
\begin{figure}[t]
\centering
\begin{tikzpicture}
\begin{groupplot}[group style={group size=1 by 2, vertical sep=0.45cm, x descriptions at=edge bottom},
   curso, width=12.4cm, xmin=0, xmax=3000]
\nextgroupplot[height=5.2cm, ylabel={metilación $\beta$}, ymin=0, ymax=1.05,
   legend style={at={(0.5,1.02)},anchor=south}, legend columns=4]
  \fill[amarillo!18] (axis cs:\capdmrini,0) rectangle (axis cs:\capdmrfin,1.05);
  \addplot[only marks, mark=*, mark size=1.1pt, azul!70] table[x=x,y=b] {figuras/cap13/dmr_a.dat};
  \addlegendentry{réplicas grupo A}
  \addplot[only marks, mark=square*, mark size=1.1pt, naranja!80] table[x=x,y=b] {figuras/cap13/dmr_b.dat};
  \addlegendentry{réplicas grupo B}
  \addplot[azul, very thick] table[x=x,y=a] {figuras/cap13/dmr_curva.dat};
  \addlegendentry{A suavizado}
  \addplot[naranja, very thick] table[x=x,y=b] {figuras/cap13/dmr_curva.dat};
  \addlegendentry{B suavizado}
  \node[etiqueta,anchor=north,text=amarillo!30!black] at (axis cs:1497,1.03) {DMR llamada};
\nextgroupplot[height=3.3cm, ylabel={$-\log_{10}p$}, xlabel={posición (pb)}, ymin=0, ymax=13]
  \fill[amarillo!18] (axis cs:\capdmrini,0) rectangle (axis cs:\capdmrfin,13);
  \addplot[ycomb, violeta, thick, mark=*, mark size=1pt] table[x=x,y=lp] {figuras/cap13/dmr_p.dat};
  \draw[rojo,dashed] (axis cs:0,2) -- (axis cs:3000,2);
  \node[etiqueta,anchor=south west,text=rojooscuro] at (axis cs:30,2.1) {$p=0{,}01$};
\end{groupplot}
\end{tikzpicture}
\caption[Detección de una región diferencialmente metilada]{Región diferencialmente metilada simulada: 55 CpG en \SI{3}{kb}, tres réplicas por grupo, cobertura media de 14 lecturas y datos beta-binomiales con $\phi=0{,}05$. \textbf{Arriba}: niveles por réplica (puntos) y curvas suavizadas con un núcleo gaussiano de 150 pb. El grupo B pierde la metilación en un bloque central, como ocurre en la hipometilación de un potenciador activado. \textbf{Abajo}: significancia de la prueba de Wald beta-binomial por CpG (\cref{eq:13-wald}). La banda amarilla es la DMR llamada como el tramo contiguo más largo de CpG con $p<0{,}01$. Los puntos individuales son muy ruidosos; la región, en cambio, es inequívoca.}
\label{fig:13-dmr}
\end{figure}

\subsection{Relojes epigenéticos y metilación sin bisulfito}

La metilación cambia con la edad de manera tan regular que puede usarse como reloj. \textcite{c13-horvath2013} entrenó un predictor multitejido con unas \num{8000} muestras de micromatrices de 82 conjuntos de datos y 51 tejidos y tipos celulares sanos: una regresión penalizada (red elástica) seleccionó 353 CpG cuya combinación lineal, tras una transformación de la edad logarítmica en la infancia y lineal en la adultez, estima la ``edad epigenética''. Esta edad es cercana a cero en células madre embrionarias e inducidas, aumenta con el número de pases en cultivo y aparece acelerada en muchos tipos de cáncer. La diferencia entre la edad epigenética y la cronológica, la \emph{aceleración}, se ha convertido en un biomarcador muy estudiado.

\paragraph{Nanoporos} La secuenciación por nanoporos (\cref{cap:06}) lee la molécula nativa, sin amplificar, y la corriente iónica depende de las bases que ocupan el poro, incluidas las modificadas. \textcite{c13-simpson2017} entrenaron un modelo oculto de Márkov con ADN metilado sintéticamente para distinguir 5mC de citosina en la señal eléctrica y obtuvieron el metiloma humano sin ningún paso especial de preparación. Sin bisulfito no hay degradación del ADN ni pérdida de complejidad, las lecturas largas permiten ``fasear'' la metilación con los alelos, y la misma corrida entrega secuencia y epigenoma.

\begin{cuidado}[Trampas de la metilación]
\begin{enumerate}
  \item \emph{C$\to$T que no es bisulfito.} Un SNP C/T en la muestra es indistinguible de una C no metilada en la hebra correspondiente; hay que filtrar posiciones polimórficas o mirar la hebra opuesta (donde el SNP aparece como G/A y el bisulfito no actúa).
  \item \emph{Composición celular.} Una diferencia de metilación entre tejidos enfermos y sanos puede reflejar solo una proporción distinta de tipos celulares. Deconvolucionar o ajustar por composición es obligatorio en muestras de sangre y tumores.
\end{enumerate}
\end{cuidado}

\begin{ideaclave}
El bisulfito convierte una marca química invisible en una diferencia de secuencia (C frente a T). La metilación de un sitio es una proporción de lecturas, y su comparación entre grupos exige un modelo que separe el ruido de muestreo (binomial) de la variación biológica (dispersión beta).
\end{ideaclave}

% ---------------------------------------------------------------------
\section{Descubrimiento de motivos regulatorios}\label{sec:13-motivos}
% ---------------------------------------------------------------------
\enlacecolab{modulo-13-epigenomica/13.3_descubrimiento_motivos.ipynb}{13.3 · Descubrimiento de motivos}

\begin{intuicion}[La contraseña escondida en treinta cartas]
Suponga que intercepta treinta cartas escritas en un alfabeto de cuatro letras, casi todas con texto sin sentido. Le han dicho que en cada una hay escondida, en algún lugar, una misma contraseña de ocho letras, pero copiada con errores. Si supiera dónde está en cada carta, bastaría alinear los fragmentos y contar letras. Si conociera la contraseña, bastaría buscarla en cada carta. No sabe ninguna de las dos cosas. La salida es circular pero funciona: adivine una contraseña aproximada, busque en cada carta el fragmento que más se le parece, reconstruya la contraseña a partir de esos fragmentos, y repita. Cada vuelta mejora un poco las dos conjeturas a la vez, hasta que dejan de cambiar.
\end{intuicion}

Un pico de ChIP-seq indica \emph{dónde} se une un factor, pero con una resolución de decenas o cientos de pares de bases. La especificidad de la unión reside en un segmento corto de 6 a 20 pb, el \emph{motivo}, que el \cref{cap:04} representó con una matriz de peso posicional (PWM) y un \ingles{sequence logo}. Allí la matriz se construía a partir de sitios conocidos. Aquí el problema es el inverso: dados cientos o miles de secuencias que \emph{probablemente} contienen el motivo en alguna posición desconocida, aprender simultáneamente la matriz y las posiciones. Es un ejemplo de manual de problema con \emph{datos incompletos}.

\subsection{El modelo generativo}

Sean $X_1,\dots,X_N$ secuencias de longitud $L$ y $W$ el ancho del motivo; cada secuencia tiene $m=L-W+1$ posiciones de inicio posibles. El modelo más sencillo, llamado OOPS (\ingles{one occurrence per sequence}), supone que cada secuencia contiene exactamente un sitio, cuyas $W$ bases se generan con la matriz $\theta$ (una distribución sobre $\{A,C,G,T\}$ por columna), mientras que el resto de las bases proviene de un fondo $p_0$. Introducimos variables \emph{ocultas} $Z_{ij}\in\{0,1\}$ que valen 1 si el sitio de la secuencia $i$ empieza en la posición $j$.

\begin{definicion}{Verosimilitud del modelo OOPS}{13-oops}
Con un inicio uniforme ($\Prob(Z_{ij}=1)=1/m$), la razón de verosimilitudes de la secuencia $X_i$ frente al modelo de fondo puro es
\begin{equation}\label{eq:13-oops}
\frac{\Prob(X_i\mid\theta)}{\Prob(X_i\mid p_0)} = \frac{1}{m}\sum_{j=1}^{m} r_{ij}(\theta),\qquad
r_{ij}(\theta)=\prod_{k=1}^{W}\frac{\theta_{k,\,x_{i,j+k-1}}}{p_0(x_{i,j+k-1})},
\end{equation}
y la log-verosimilitud del conjunto, en bits, es $\mathrm{LLR}(\theta)=\sum_i\log_2\bigl(\frac1m\sum_j r_{ij}(\theta)\bigr)$.
\end{definicion}

\begin{simbolos}
\simbolo{$x_{i,j}$}{Base en la posición $j$ de la secuencia $i$.}
\simbolo{$\theta_{k,b}$}{Probabilidad de la base $b$ en la columna $k$ del motivo.}
\simbolo{$p_0(b)$}{Frecuencia de fondo de la base $b$.}
\simbolo{$r_{ij}$}{Razón de verosimilitudes de que el sitio esté en $j$ frente a que sea fondo; $\log_2 r_{ij}$ es la puntuación PWM del \cref{cap:04}.}
\simbolo{$m$}{Número de posiciones de inicio posibles, $L-W+1$.}
\end{simbolos}

Las bases fuera del sitio se cancelan en el cociente, así que solo quedan las $W$ columnas del candidato. Si conociéramos las $Z_{ij}$, estimar $\theta$ sería contar; si conociéramos $\theta$, localizar los sitios sería escanear con la PWM. El algoritmo EM convierte esta circularidad en un procedimiento con garantías.

\subsection{El algoritmo EM y MEME}

\textcite{c13-dempster1977} formularon el algoritmo de esperanza-maximización (EM) en general para problemas con datos faltantes. Aplicado a motivos alterna dos pasos:

\begin{teorema}{EM para el modelo OOPS}{13-em}
Dada una estimación $\theta^{(t)}$:
\begin{align}
\text{Paso E:}\quad & Z_{ij}^{(t)} = \Prob\bigl(Z_{ij}=1\mid X_i,\theta^{(t)}\bigr) = \frac{r_{ij}(\theta^{(t)})}{\sum_{j'=1}^{m} r_{ij'}(\theta^{(t)})},\label{eq:13-estep}\\[2pt]
\text{Paso M:}\quad & \theta^{(t+1)}_{k,b} = \frac{\sum_{i}\sum_{j} Z_{ij}^{(t)}\,\mathbb 1\bigl[x_{i,j+k-1}=b\bigr] + \beta_b}{\sum_{b'}\Bigl(\sum_{i}\sum_{j} Z_{ij}^{(t)}\,\mathbb 1\bigl[x_{i,j+k-1}=b'\bigr] + \beta_{b'}\Bigr)} .\label{eq:13-mstep}
\end{align}
Con $\beta_b=0$, cada iteración no disminuye la verosimilitud: $\mathrm{LLR}(\theta^{(t+1)})\ge\mathrm{LLR}(\theta^{(t)})$.
\end{teorema}

\begin{simbolos}
\simbolo{$Z_{ij}^{(t)}$}{Probabilidad posterior de que el sitio de la secuencia $i$ empiece en $j$; es un ``conteo fraccionario''.}
\simbolo{$\beta_b$}{Seudoconteos (a priori de Dirichlet) que evitan probabilidades nulas.}
\simbolo{$\mathbb 1[\cdot]$}{Función indicadora.}
\end{simbolos}

El paso E reparte cada secuencia entre sus posiciones posibles en proporción a cuánto se parecen al motivo actual: en vez de decidir dónde está el sitio, reparte ``votos''. El paso M construye una nueva PWM con esos votos fraccionarios. La garantía de monotonía viene de la desigualdad de Jensen: para cualquier distribución $q$ sobre las variables ocultas, $\log\Prob(X\mid\theta)\ge\sum_Z q(Z)\log\frac{\Prob(X,Z\mid\theta)}{q(Z)}$, con igualdad cuando $q$ es la posterior. El paso E hace la cota exacta en $\theta^{(t)}$; el paso M la maximiza en $\theta$; por tanto la verosimilitud en $\theta^{(t+1)}$ es al menos la cota, que es al menos la verosimilitud en $\theta^{(t)}$. Lo que EM \emph{no} garantiza es el máximo global.

\begin{ejemplo}{Un paso E}{13-estep}
En una colección simulada de 30 secuencias de 80 pb con un motivo de tipo AP-1 implantado (consenso \secuencia{ATGASTCA}, $W=8$; hay $m=73$ posiciones por secuencia), la secuencia 18 contiene el sitio \secuencia{ATGACTCA}, idéntico al consenso. Con la matriz verdadera y fondo uniforme, su puntuación es
\[
\log_2 r = 1{,}26+1{,}77+1{,}77+1{,}82+0{,}91+1{,}77+1{,}77+1{,}82 \approx 12{,}87\ \text{bits},
\]
es decir, $r\approx 7500$: el sitio es unas siete mil quinientas veces más probable bajo el motivo que bajo el fondo. Al normalizar sobre las 73 posiciones, el paso E le asigna $Z=0{,}9976$ y reparte el $0{,}24\,\%$ restante entre las otras 72. Pero no todos los sitios son tan claros: la mediana de las puntuaciones de los 30 sitios reales es de 8,65 bits, y en las secuencias con sitios degenerados el paso E reparte el voto entre varias posiciones. Cambiar una sola base (\secuencia{ATGACTTA}) baja la puntuación a 8,78 bits.
\end{ejemplo}

\begin{codigo}[em\_oops.py]
import numpy as np

def em_oops(seqs, theta, iters=30, beta=0.1):
    """seqs: matriz N x L de enteros 0-3; theta: W x 4 inicial."""
    N, L = seqs.shape
    W = theta.shape[0]
    win = np.stack([seqs[:, j:j + W] for j in range(L - W + 1)], 1)
    for _ in range(iters):
        lr = np.log2(theta[np.arange(W), win]).sum(-1) + 2 * W
        Z = np.exp2(lr - lr.max(1, keepdims=True))     # paso E
        Z /= Z.sum(1, keepdims=True)
        cnt = np.stack([(Z[..., None] * (win == b)).sum((0, 1))
                        for b in range(4)], 1)         # paso M
        theta = (cnt + beta) / (cnt + beta).sum(1, keepdims=True)
    return theta, Z
\end{codigo}

\paragraph{MEME} El programa MEME, presentado por Bailey y Elkan en 1994 y hoy núcleo de la MEME Suite \parencite{c13-bailey2009}, es la implementación de referencia de esta idea. Añade tres ingredientes decisivos. Primero, modelos más realistas que OOPS: ZOOPS (\ingles{zero or one occurrence per sequence}), que admite secuencias sin sitio, y TCM (\ingles{two-component mixture}), que admite cualquier número. Segundo, una estrategia de \emph{inicialización}: en lugar de partir de una matriz aleatoria, prueba como semillas subcadenas reales de los datos, corre una iteración corta desde cada una y lanza EM completo solo desde las más prometedoras. Tercero, un criterio estadístico (el valor $E$ del motivo) para decidir cuántos motivos reportar y de qué ancho.

La \cref{fig:13-em}\,(a) muestra por qué la inicialización importa. Seis corridas de EM sobre el mismo conjunto, cinco con semillas aleatorias y una con una subcadena de la secuencia, terminan en tres destinos. Dos alcanzan la verosimilitud máxima encontrada, $\mathrm{LLR}=96{,}4$ bits, con el consenso \secuencia{ATGAGTCA} (la S del motivo admite C o G). Otras quedan atrapadas en versiones \emph{desfasadas} del motivo (\secuencia{TGAGTCAT}, \secuencia{CATGACTC}), que capturan solo siete de las ocho columnas, o en un óptimo local pobre ($\mathrm{LLR}=56{,}5$). La curva de la corrida ganadora tiene forma de S: durante cinco iteraciones apenas mejora, hasta que un puñado de sitios reales ``tira'' de la matriz y el resto la sigue en cascada. En esa corrida, 26 de los 30 sitios implantados se recuperan en la posición exacta, y la matriz final tiene 8,75 bits de información, \emph{más} que la verdadera (8,11): con solo 30 sitios, la máxima verosimilitud sobreajusta y hace el motivo más nítido de lo que es. Por la misma razón, la $\mathrm{LLR}$ de la solución (96,4 bits) supera a la de la matriz verdadera (76,1).

\begin{figure}[t]
\centering
\begin{tikzpicture}
\begin{groupplot}[group style={group size=2 by 1, horizontal sep=1.35cm},
   curso, height=5.4cm, ymin=-10, ymax=108]
\nextgroupplot[width=6.5cm, title={(a) EM desde seis semillas},
   xlabel={iteración}, ylabel={LLR (bits)}, xmin=0, xmax=30]
  \addplot[gris, thick] table[x=it,y=r0] {figuras/cap13/em_ll.dat};
  \addplot[aqua, thick] table[x=it,y=r1] {figuras/cap13/em_ll.dat};
  \addplot[magenta, thick] table[x=it,y=r3] {figuras/cap13/em_ll.dat};
  \addplot[violeta, thick] table[x=it,y=r4] {figuras/cap13/em_ll.dat};
  \addplot[naranja, very thick, dashed] table[x=it,y=r5] {figuras/cap13/em_ll.dat};
  \addplot[azul, very thick] table[x=it,y=r2] {figuras/cap13/em_ll.dat};
  \draw[rojo,dashed] (axis cs:0,76.1) -- (axis cs:30,76.1);
  \node[etiqueta,anchor=south west,text=rojooscuro,fill=fondo,inner sep=1pt] at (axis cs:0.5,79) {matriz verdadera};
\nextgroupplot[width=6.5cm, title={(b) Gibbs: cuatro cadenas},
   xlabel={barrido}, ylabel={LLR (bits)}, xmin=0, xmax=60]
  \addplot[gris, thick] table[x=it,y=c0] {figuras/cap13/gibbs_ll.dat};
  \addplot[aqua, thick] table[x=it,y=c1] {figuras/cap13/gibbs_ll.dat};
  \addplot[magenta, thick] table[x=it,y=c2] {figuras/cap13/gibbs_ll.dat};
  \addplot[azul, very thick] table[x=it,y=c3] {figuras/cap13/gibbs_ll.dat};
  \draw[rojo,dashed] (axis cs:0,76.1) -- (axis cs:60,76.1);
  \node[etiqueta,anchor=south,text=azulprofundo] at (axis cs:40,97) {\texttt{ATGACTCA}};
\end{groupplot}
\end{tikzpicture}
\caption[Convergencia de EM y de Gibbs]{Descubrimiento de un motivo de 8 pb implantado en 30 secuencias de 80 pb. (a) Log-verosimilitud (\cref{eq:13-oops}) a lo largo de las iteraciones de EM desde seis puntos de partida: cinco semillas aleatorias y una subcadena de los datos (naranja discontinua), como hace MEME. Todas las curvas son crecientes, como garantiza el \cref{teo:13-em}, pero terminan en óptimos distintos; la mejor corrida con semilla aleatoria (azul) muestra la típica curva en S. (b) Cuatro cadenas del muestreador de Gibbs: la trayectoria es ruidosa porque cada paso es aleatorio. Una cadena (azul) encuentra el motivo en seis barridos; dos quedan desfasadas dos posiciones y una no encuentra nada. La línea roja es la verosimilitud de la matriz verdadera.}
\label{fig:13-em}
\end{figure}

\begin{figure}[t]
\centering
\begin{tabular}{@{}ccc@{}}
\begin{tikzpicture}[x=0.42cm,y=0.62cm]
\draw[base] (0.4,0) -- (8.60,0);
\draw[base] (0.4,0) -- (0.4,2);
\draw[base] (0.4,0) -- (0.3,0); \node[font=\sffamily\tiny,text=gris,anchor=east] at (0.3,0) {0};
\draw[base] (0.4,1) -- (0.3,1); \node[font=\sffamily\tiny,text=gris,anchor=east] at (0.3,1) {1};
\draw[base] (0.4,2) -- (0.3,2); \node[font=\sffamily\tiny,text=gris,anchor=east] at (0.3,2) {2};
\node[anchor=south west,inner sep=0pt,text=nucT] at (0.560,0.0000) {\resizebox{0.370cm}{0.0251cm}{\sffamily\bfseries T}};
\node[anchor=south west,inner sep=0pt,text=nucC] at (0.560,0.0405) {\resizebox{0.370cm}{0.0376cm}{\sffamily\bfseries C}};
\node[anchor=south west,inner sep=0pt,text=nucG] at (0.560,0.1011) {\resizebox{0.370cm}{0.0376cm}{\sffamily\bfseries G}};
\node[anchor=south west,inner sep=0pt,text=nucA] at (0.560,0.1618) {\resizebox{0.370cm}{0.1505cm}{\sffamily\bfseries A}};
\node[font=\sffamily\tiny,text=gris] at (1,-0.2) {1};
\node[anchor=south west,inner sep=0pt,text=nucA] at (1.560,0.0000) {\resizebox{0.370cm}{0.0357cm}{\sffamily\bfseries A}};
\node[anchor=south west,inner sep=0pt,text=nucC] at (1.560,0.0576) {\resizebox{0.370cm}{0.0357cm}{\sffamily\bfseries C}};
\node[anchor=south west,inner sep=0pt,text=nucG] at (1.560,0.1152) {\resizebox{0.370cm}{0.0357cm}{\sffamily\bfseries G}};
\node[anchor=south west,inner sep=0pt,text=nucT] at (1.560,0.1729) {\resizebox{0.370cm}{0.6073cm}{\sffamily\bfseries T}};
\node[font=\sffamily\tiny,text=gris] at (2,-0.2) {2};
\node[anchor=south west,inner sep=0pt,text=nucA] at (2.560,0.0000) {\resizebox{0.370cm}{0.0357cm}{\sffamily\bfseries A}};
\node[anchor=south west,inner sep=0pt,text=nucC] at (2.560,0.0576) {\resizebox{0.370cm}{0.0357cm}{\sffamily\bfseries C}};
\node[anchor=south west,inner sep=0pt,text=nucT] at (2.560,0.1152) {\resizebox{0.370cm}{0.0357cm}{\sffamily\bfseries T}};
\node[anchor=south west,inner sep=0pt,text=nucG] at (2.560,0.1729) {\resizebox{0.370cm}{0.6073cm}{\sffamily\bfseries G}};
\node[font=\sffamily\tiny,text=gris] at (3,-0.2) {3};
\node[anchor=south west,inner sep=0pt,text=nucC] at (3.560,0.0000) {\resizebox{0.370cm}{0.0318cm}{\sffamily\bfseries C}};
\node[anchor=south west,inner sep=0pt,text=nucG] at (3.560,0.0512) {\resizebox{0.370cm}{0.0318cm}{\sffamily\bfseries G}};
\node[anchor=south west,inner sep=0pt,text=nucT] at (3.560,0.1024) {\resizebox{0.370cm}{0.0318cm}{\sffamily\bfseries T}};
\node[anchor=south west,inner sep=0pt,text=nucA] at (3.560,0.1537) {\resizebox{0.370cm}{0.6986cm}{\sffamily\bfseries A}};
\node[font=\sffamily\tiny,text=gris] at (4,-0.2) {4};
\node[anchor=south west,inner sep=0pt,text=nucA] at (4.560,0.0000) {\resizebox{0.370cm}{0.0165cm}{\sffamily\bfseries A}};
\node[anchor=south west,inner sep=0pt,text=nucT] at (4.560,0.0266) {\resizebox{0.370cm}{0.0165cm}{\sffamily\bfseries T}};
\node[anchor=south west,inner sep=0pt,text=nucG] at (4.560,0.0532) {\resizebox{0.370cm}{0.1419cm}{\sffamily\bfseries G}};
\node[anchor=south west,inner sep=0pt,text=nucC] at (4.560,0.2821) {\resizebox{0.370cm}{0.1551cm}{\sffamily\bfseries C}};
\node[font=\sffamily\tiny,text=gris] at (5,-0.2) {5};
\node[anchor=south west,inner sep=0pt,text=nucA] at (5.560,0.0000) {\resizebox{0.370cm}{0.0357cm}{\sffamily\bfseries A}};
\node[anchor=south west,inner sep=0pt,text=nucC] at (5.560,0.0576) {\resizebox{0.370cm}{0.0357cm}{\sffamily\bfseries C}};
\node[anchor=south west,inner sep=0pt,text=nucG] at (5.560,0.1152) {\resizebox{0.370cm}{0.0357cm}{\sffamily\bfseries G}};
\node[anchor=south west,inner sep=0pt,text=nucT] at (5.560,0.1729) {\resizebox{0.370cm}{0.6073cm}{\sffamily\bfseries T}};
\node[font=\sffamily\tiny,text=gris] at (6,-0.2) {6};
\node[anchor=south west,inner sep=0pt,text=nucA] at (6.560,0.0000) {\resizebox{0.370cm}{0.0357cm}{\sffamily\bfseries A}};
\node[anchor=south west,inner sep=0pt,text=nucG] at (6.560,0.0576) {\resizebox{0.370cm}{0.0357cm}{\sffamily\bfseries G}};
\node[anchor=south west,inner sep=0pt,text=nucT] at (6.560,0.1152) {\resizebox{0.370cm}{0.0357cm}{\sffamily\bfseries T}};
\node[anchor=south west,inner sep=0pt,text=nucC] at (6.560,0.1729) {\resizebox{0.370cm}{0.6073cm}{\sffamily\bfseries C}};
\node[font=\sffamily\tiny,text=gris] at (7,-0.2) {7};
\node[anchor=south west,inner sep=0pt,text=nucC] at (7.560,0.0000) {\resizebox{0.370cm}{0.0318cm}{\sffamily\bfseries C}};
\node[anchor=south west,inner sep=0pt,text=nucG] at (7.560,0.0512) {\resizebox{0.370cm}{0.0318cm}{\sffamily\bfseries G}};
\node[anchor=south west,inner sep=0pt,text=nucT] at (7.560,0.1024) {\resizebox{0.370cm}{0.0318cm}{\sffamily\bfseries T}};
\node[anchor=south west,inner sep=0pt,text=nucA] at (7.560,0.1537) {\resizebox{0.370cm}{0.6986cm}{\sffamily\bfseries A}};
\node[font=\sffamily\tiny,text=gris] at (8,-0.2) {8};
\node[font=\sffamily\scriptsize\bfseries,text=tinta,anchor=south] at (4.50,2.05) {motivo implantado};
\end{tikzpicture} & \begin{tikzpicture}[x=0.42cm,y=0.62cm]
\draw[base] (0.4,0) -- (8.60,0);
\draw[base] (0.4,0) -- (0.4,2);
\draw[base] (0.4,0) -- (0.3,0); \node[font=\sffamily\tiny,text=gris,anchor=east] at (0.3,0) {0};
\draw[base] (0.4,1) -- (0.3,1); \node[font=\sffamily\tiny,text=gris,anchor=east] at (0.3,1) {1};
\draw[base] (0.4,2) -- (0.3,2); \node[font=\sffamily\tiny,text=gris,anchor=east] at (0.3,2) {2};
\node[anchor=south west,inner sep=0pt,text=nucA] at (0.560,0.0000) {\resizebox{0.370cm}{0.0214cm}{\sffamily\bfseries A}};
\node[anchor=south west,inner sep=0pt,text=nucC] at (0.560,0.0346) {\resizebox{0.370cm}{0.0214cm}{\sffamily\bfseries C}};
\node[anchor=south west,inner sep=0pt,text=nucG] at (0.560,0.0692) {\resizebox{0.370cm}{0.0214cm}{\sffamily\bfseries G}};
\node[anchor=south west,inner sep=0pt,text=nucT] at (0.560,0.1038) {\resizebox{0.370cm}{0.0643cm}{\sffamily\bfseries T}};
\node[font=\sffamily\tiny,text=gris] at (1,-0.2) {1};
\node[anchor=south west,inner sep=0pt,text=nucA] at (1.560,0.0000) {\resizebox{0.370cm}{0.0214cm}{\sffamily\bfseries A}};
\node[anchor=south west,inner sep=0pt,text=nucC] at (1.560,0.0346) {\resizebox{0.370cm}{0.0214cm}{\sffamily\bfseries C}};
\node[anchor=south west,inner sep=0pt,text=nucG] at (1.560,0.0692) {\resizebox{0.370cm}{0.0214cm}{\sffamily\bfseries G}};
\node[anchor=south west,inner sep=0pt,text=nucT] at (1.560,0.1038) {\resizebox{0.370cm}{0.0643cm}{\sffamily\bfseries T}};
\node[font=\sffamily\tiny,text=gris] at (2,-0.2) {2};
\node[anchor=south west,inner sep=0pt,text=nucA] at (2.560,0.0000) {\resizebox{0.370cm}{0.0214cm}{\sffamily\bfseries A}};
\node[anchor=south west,inner sep=0pt,text=nucG] at (2.560,0.0346) {\resizebox{0.370cm}{0.0214cm}{\sffamily\bfseries G}};
\node[anchor=south west,inner sep=0pt,text=nucT] at (2.560,0.0692) {\resizebox{0.370cm}{0.0214cm}{\sffamily\bfseries T}};
\node[anchor=south west,inner sep=0pt,text=nucC] at (2.560,0.1038) {\resizebox{0.370cm}{0.0643cm}{\sffamily\bfseries C}};
\node[font=\sffamily\tiny,text=gris] at (3,-0.2) {3};
\node[anchor=south west,inner sep=0pt,text=nucA] at (3.560,0.0000) {\resizebox{0.370cm}{0.0214cm}{\sffamily\bfseries A}};
\node[anchor=south west,inner sep=0pt,text=nucG] at (3.560,0.0346) {\resizebox{0.370cm}{0.0214cm}{\sffamily\bfseries G}};
\node[anchor=south west,inner sep=0pt,text=nucT] at (3.560,0.0692) {\resizebox{0.370cm}{0.0214cm}{\sffamily\bfseries T}};
\node[anchor=south west,inner sep=0pt,text=nucC] at (3.560,0.1038) {\resizebox{0.370cm}{0.0643cm}{\sffamily\bfseries C}};
\node[font=\sffamily\tiny,text=gris] at (4,-0.2) {4};
\node[anchor=south west,inner sep=0pt,text=nucA] at (4.560,0.0000) {\resizebox{0.370cm}{0.0214cm}{\sffamily\bfseries A}};
\node[anchor=south west,inner sep=0pt,text=nucG] at (4.560,0.0346) {\resizebox{0.370cm}{0.0214cm}{\sffamily\bfseries G}};
\node[anchor=south west,inner sep=0pt,text=nucT] at (4.560,0.0692) {\resizebox{0.370cm}{0.0214cm}{\sffamily\bfseries T}};
\node[anchor=south west,inner sep=0pt,text=nucC] at (4.560,0.1038) {\resizebox{0.370cm}{0.0643cm}{\sffamily\bfseries C}};
\node[font=\sffamily\tiny,text=gris] at (5,-0.2) {5};
\node[anchor=south west,inner sep=0pt,text=nucA] at (5.560,0.0000) {\resizebox{0.370cm}{0.0214cm}{\sffamily\bfseries A}};
\node[anchor=south west,inner sep=0pt,text=nucC] at (5.560,0.0346) {\resizebox{0.370cm}{0.0214cm}{\sffamily\bfseries C}};
\node[anchor=south west,inner sep=0pt,text=nucG] at (5.560,0.0692) {\resizebox{0.370cm}{0.0214cm}{\sffamily\bfseries G}};
\node[anchor=south west,inner sep=0pt,text=nucT] at (5.560,0.1038) {\resizebox{0.370cm}{0.0643cm}{\sffamily\bfseries T}};
\node[font=\sffamily\tiny,text=gris] at (6,-0.2) {6};
\node[anchor=south west,inner sep=0pt,text=nucA] at (6.560,0.0000) {\resizebox{0.370cm}{0.0214cm}{\sffamily\bfseries A}};
\node[anchor=south west,inner sep=0pt,text=nucC] at (6.560,0.0346) {\resizebox{0.370cm}{0.0214cm}{\sffamily\bfseries C}};
\node[anchor=south west,inner sep=0pt,text=nucG] at (6.560,0.0692) {\resizebox{0.370cm}{0.0214cm}{\sffamily\bfseries G}};
\node[anchor=south west,inner sep=0pt,text=nucT] at (6.560,0.1038) {\resizebox{0.370cm}{0.0643cm}{\sffamily\bfseries T}};
\node[font=\sffamily\tiny,text=gris] at (7,-0.2) {7};
\node[anchor=south west,inner sep=0pt,text=nucC] at (7.560,0.0000) {\resizebox{0.370cm}{0.0214cm}{\sffamily\bfseries C}};
\node[anchor=south west,inner sep=0pt,text=nucG] at (7.560,0.0346) {\resizebox{0.370cm}{0.0214cm}{\sffamily\bfseries G}};
\node[anchor=south west,inner sep=0pt,text=nucT] at (7.560,0.0692) {\resizebox{0.370cm}{0.0214cm}{\sffamily\bfseries T}};
\node[anchor=south west,inner sep=0pt,text=nucA] at (7.560,0.1038) {\resizebox{0.370cm}{0.0643cm}{\sffamily\bfseries A}};
\node[font=\sffamily\tiny,text=gris] at (8,-0.2) {8};
\node[font=\sffamily\scriptsize\bfseries,text=tinta,anchor=south] at (4.50,2.05) {semilla (iteración 0)};
\end{tikzpicture} & \begin{tikzpicture}[x=0.42cm,y=0.62cm]
\draw[base] (0.4,0) -- (8.60,0);
\draw[base] (0.4,0) -- (0.4,2);
\draw[base] (0.4,0) -- (0.3,0); \node[font=\sffamily\tiny,text=gris,anchor=east] at (0.3,0) {0};
\draw[base] (0.4,1) -- (0.3,1); \node[font=\sffamily\tiny,text=gris,anchor=east] at (0.3,1) {1};
\draw[base] (0.4,2) -- (0.3,2); \node[font=\sffamily\tiny,text=gris,anchor=east] at (0.3,2) {2};
\node[anchor=south west,inner sep=0pt,text=nucC] at (0.560,0.0000) {\resizebox{0.370cm}{0.0153cm}{\sffamily\bfseries C}};
\node[anchor=south west,inner sep=0pt,text=nucG] at (0.560,0.0247) {\resizebox{0.370cm}{0.0286cm}{\sffamily\bfseries G}};
\node[anchor=south west,inner sep=0pt,text=nucA] at (0.560,0.0708) {\resizebox{0.370cm}{0.0387cm}{\sffamily\bfseries A}};
\node[anchor=south west,inner sep=0pt,text=nucT] at (0.560,0.1332) {\resizebox{0.370cm}{0.0870cm}{\sffamily\bfseries T}};
\node[font=\sffamily\tiny,text=gris] at (1,-0.2) {1};
\node[anchor=south west,inner sep=0pt,text=nucG] at (1.560,0.0000) {\resizebox{0.370cm}{0.0218cm}{\sffamily\bfseries G}};
\node[anchor=south west,inner sep=0pt,text=nucC] at (1.560,0.0352) {\resizebox{0.370cm}{0.0243cm}{\sffamily\bfseries C}};
\node[anchor=south west,inner sep=0pt,text=nucA] at (1.560,0.0744) {\resizebox{0.370cm}{0.0472cm}{\sffamily\bfseries A}};
\node[anchor=south west,inner sep=0pt,text=nucT] at (1.560,0.1506) {\resizebox{0.370cm}{0.1137cm}{\sffamily\bfseries T}};
\node[font=\sffamily\tiny,text=gris] at (2,-0.2) {2};
\node[anchor=south west,inner sep=0pt,text=nucT] at (2.560,0.0000) {\resizebox{0.370cm}{0.0118cm}{\sffamily\bfseries T}};
\node[anchor=south west,inner sep=0pt,text=nucA] at (2.560,0.0190) {\resizebox{0.370cm}{0.0247cm}{\sffamily\bfseries A}};
\node[anchor=south west,inner sep=0pt,text=nucG] at (2.560,0.0588) {\resizebox{0.370cm}{0.0700cm}{\sffamily\bfseries G}};
\node[anchor=south west,inner sep=0pt,text=nucC] at (2.560,0.1717) {\resizebox{0.370cm}{0.0713cm}{\sffamily\bfseries C}};
\node[font=\sffamily\tiny,text=gris] at (3,-0.2) {3};
\node[anchor=south west,inner sep=0pt,text=nucG] at (3.560,0.0000) {\resizebox{0.370cm}{0.0119cm}{\sffamily\bfseries G}};
\node[anchor=south west,inner sep=0pt,text=nucT] at (3.560,0.0191) {\resizebox{0.370cm}{0.0238cm}{\sffamily\bfseries T}};
\node[anchor=south west,inner sep=0pt,text=nucA] at (3.560,0.0575) {\resizebox{0.370cm}{0.0474cm}{\sffamily\bfseries A}};
\node[anchor=south west,inner sep=0pt,text=nucC] at (3.560,0.1339) {\resizebox{0.370cm}{0.0657cm}{\sffamily\bfseries C}};
\node[font=\sffamily\tiny,text=gris] at (4,-0.2) {4};
\node[anchor=south west,inner sep=0pt,text=nucT] at (4.560,0.0000) {\resizebox{0.370cm}{0.0136cm}{\sffamily\bfseries T}};
\node[anchor=south west,inner sep=0pt,text=nucA] at (4.560,0.0219) {\resizebox{0.370cm}{0.0263cm}{\sffamily\bfseries A}};
\node[anchor=south west,inner sep=0pt,text=nucC] at (4.560,0.0644) {\resizebox{0.370cm}{0.0577cm}{\sffamily\bfseries C}};
\node[anchor=south west,inner sep=0pt,text=nucG] at (4.560,0.1575) {\resizebox{0.370cm}{0.0934cm}{\sffamily\bfseries G}};
\node[font=\sffamily\tiny,text=gris] at (5,-0.2) {5};
\node[anchor=south west,inner sep=0pt,text=nucC] at (5.560,0.0000) {\resizebox{0.370cm}{0.0360cm}{\sffamily\bfseries C}};
\node[anchor=south west,inner sep=0pt,text=nucA] at (5.560,0.0581) {\resizebox{0.370cm}{0.0365cm}{\sffamily\bfseries A}};
\node[anchor=south west,inner sep=0pt,text=nucG] at (5.560,0.1169) {\resizebox{0.370cm}{0.0493cm}{\sffamily\bfseries G}};
\node[anchor=south west,inner sep=0pt,text=nucT] at (5.560,0.1965) {\resizebox{0.370cm}{0.3022cm}{\sffamily\bfseries T}};
\node[font=\sffamily\tiny,text=gris] at (6,-0.2) {6};
\node[anchor=south west,inner sep=0pt,text=nucA] at (6.560,0.0000) {\resizebox{0.370cm}{0.0126cm}{\sffamily\bfseries A}};
\node[anchor=south west,inner sep=0pt,text=nucG] at (6.560,0.0202) {\resizebox{0.370cm}{0.0241cm}{\sffamily\bfseries G}};
\node[anchor=south west,inner sep=0pt,text=nucT] at (6.560,0.0591) {\resizebox{0.370cm}{0.0500cm}{\sffamily\bfseries T}};
\node[anchor=south west,inner sep=0pt,text=nucC] at (6.560,0.1397) {\resizebox{0.370cm}{0.0723cm}{\sffamily\bfseries C}};
\node[font=\sffamily\tiny,text=gris] at (7,-0.2) {7};
\node[anchor=south west,inner sep=0pt,text=nucG] at (7.560,0.0000) {\resizebox{0.370cm}{0.0393cm}{\sffamily\bfseries G}};
\node[anchor=south west,inner sep=0pt,text=nucC] at (7.560,0.0634) {\resizebox{0.370cm}{0.0393cm}{\sffamily\bfseries C}};
\node[anchor=south west,inner sep=0pt,text=nucT] at (7.560,0.1269) {\resizebox{0.370cm}{0.0442cm}{\sffamily\bfseries T}};
\node[anchor=south west,inner sep=0pt,text=nucA] at (7.560,0.1982) {\resizebox{0.370cm}{0.3773cm}{\sffamily\bfseries A}};
\node[font=\sffamily\tiny,text=gris] at (8,-0.2) {8};
\node[font=\sffamily\scriptsize\bfseries,text=tinta,anchor=south] at (4.50,2.05) {EM, iteración 6};
\end{tikzpicture}\\[4pt]
\begin{tikzpicture}[x=0.42cm,y=0.62cm]
\draw[base] (0.4,0) -- (8.60,0);
\draw[base] (0.4,0) -- (0.4,2);
\draw[base] (0.4,0) -- (0.3,0); \node[font=\sffamily\tiny,text=gris,anchor=east] at (0.3,0) {0};
\draw[base] (0.4,1) -- (0.3,1); \node[font=\sffamily\tiny,text=gris,anchor=east] at (0.3,1) {1};
\draw[base] (0.4,2) -- (0.3,2); \node[font=\sffamily\tiny,text=gris,anchor=east] at (0.3,2) {2};
\node[anchor=south west,inner sep=0pt,text=nucC] at (0.560,0.0000) {\resizebox{0.370cm}{0.0121cm}{\sffamily\bfseries C}};
\node[anchor=south west,inner sep=0pt,text=nucG] at (0.560,0.0195) {\resizebox{0.370cm}{0.0266cm}{\sffamily\bfseries G}};
\node[anchor=south west,inner sep=0pt,text=nucT] at (0.560,0.0623) {\resizebox{0.370cm}{0.0787cm}{\sffamily\bfseries T}};
\node[anchor=south west,inner sep=0pt,text=nucA] at (0.560,0.1893) {\resizebox{0.370cm}{0.0985cm}{\sffamily\bfseries A}};
\node[font=\sffamily\tiny,text=gris] at (1,-0.2) {1};
\node[anchor=south west,inner sep=0pt,text=nucC] at (1.560,0.0000) {\resizebox{0.370cm}{0.0212cm}{\sffamily\bfseries C}};
\node[anchor=south west,inner sep=0pt,text=nucG] at (1.560,0.0341) {\resizebox{0.370cm}{0.0239cm}{\sffamily\bfseries G}};
\node[anchor=south west,inner sep=0pt,text=nucA] at (1.560,0.0727) {\resizebox{0.370cm}{0.0805cm}{\sffamily\bfseries A}};
\node[anchor=south west,inner sep=0pt,text=nucT] at (1.560,0.2025) {\resizebox{0.370cm}{0.5121cm}{\sffamily\bfseries T}};
\node[font=\sffamily\tiny,text=gris] at (2,-0.2) {2};
\node[anchor=south west,inner sep=0pt,text=nucA] at (2.560,0.0000) {\resizebox{0.370cm}{0.0164cm}{\sffamily\bfseries A}};
\node[anchor=south west,inner sep=0pt,text=nucT] at (2.560,0.0264) {\resizebox{0.370cm}{0.0268cm}{\sffamily\bfseries T}};
\node[anchor=south west,inner sep=0pt,text=nucC] at (2.560,0.0696) {\resizebox{0.370cm}{0.0997cm}{\sffamily\bfseries C}};
\node[anchor=south west,inner sep=0pt,text=nucG] at (2.560,0.2304) {\resizebox{0.370cm}{0.2920cm}{\sffamily\bfseries G}};
\node[font=\sffamily\tiny,text=gris] at (3,-0.2) {3};
\node[anchor=south west,inner sep=0pt,text=nucG] at (3.560,0.0000) {\resizebox{0.370cm}{0.0144cm}{\sffamily\bfseries G}};
\node[anchor=south west,inner sep=0pt,text=nucT] at (3.560,0.0232) {\resizebox{0.370cm}{0.0341cm}{\sffamily\bfseries T}};
\node[anchor=south west,inner sep=0pt,text=nucC] at (3.560,0.0782) {\resizebox{0.370cm}{0.0863cm}{\sffamily\bfseries C}};
\node[anchor=south west,inner sep=0pt,text=nucA] at (3.560,0.2174) {\resizebox{0.370cm}{0.2379cm}{\sffamily\bfseries A}};
\node[font=\sffamily\tiny,text=gris] at (4,-0.2) {4};
\node[anchor=south west,inner sep=0pt,text=nucA] at (4.560,0.0087) {\resizebox{0.370cm}{0.0221cm}{\sffamily\bfseries A}};
\node[anchor=south west,inner sep=0pt,text=nucC] at (4.560,0.0443) {\resizebox{0.370cm}{0.1606cm}{\sffamily\bfseries C}};
\node[anchor=south west,inner sep=0pt,text=nucG] at (4.560,0.3034) {\resizebox{0.370cm}{0.2605cm}{\sffamily\bfseries G}};
\node[font=\sffamily\tiny,text=gris] at (5,-0.2) {5};
\node[anchor=south west,inner sep=0pt,text=nucC] at (5.560,0.0000) {\resizebox{0.370cm}{0.0134cm}{\sffamily\bfseries C}};
\node[anchor=south west,inner sep=0pt,text=nucA] at (5.560,0.0216) {\resizebox{0.370cm}{0.0380cm}{\sffamily\bfseries A}};
\node[anchor=south west,inner sep=0pt,text=nucG] at (5.560,0.0829) {\resizebox{0.370cm}{0.0643cm}{\sffamily\bfseries G}};
\node[anchor=south west,inner sep=0pt,text=nucT] at (5.560,0.1866) {\resizebox{0.370cm}{0.5835cm}{\sffamily\bfseries T}};
\node[font=\sffamily\tiny,text=gris] at (6,-0.2) {6};
\node[anchor=south west,inner sep=0pt,text=nucG] at (6.560,0.0000) {\resizebox{0.370cm}{0.0312cm}{\sffamily\bfseries G}};
\node[anchor=south west,inner sep=0pt,text=nucT] at (6.560,0.0503) {\resizebox{0.370cm}{0.0441cm}{\sffamily\bfseries T}};
\node[anchor=south west,inner sep=0pt,text=nucA] at (6.560,0.1215) {\resizebox{0.370cm}{0.0458cm}{\sffamily\bfseries A}};
\node[anchor=south west,inner sep=0pt,text=nucC] at (6.560,0.1953) {\resizebox{0.370cm}{0.2871cm}{\sffamily\bfseries C}};
\node[font=\sffamily\tiny,text=gris] at (7,-0.2) {7};
\node[anchor=south west,inner sep=0pt,text=nucC] at (7.560,0.0000) {\resizebox{0.370cm}{0.0240cm}{\sffamily\bfseries C}};
\node[anchor=south west,inner sep=0pt,text=nucG] at (7.560,0.0386) {\resizebox{0.370cm}{0.0304cm}{\sffamily\bfseries G}};
\node[anchor=south west,inner sep=0pt,text=nucT] at (7.560,0.0877) {\resizebox{0.370cm}{0.0501cm}{\sffamily\bfseries T}};
\node[anchor=south west,inner sep=0pt,text=nucA] at (7.560,0.1686) {\resizebox{0.370cm}{0.6418cm}{\sffamily\bfseries A}};
\node[font=\sffamily\tiny,text=gris] at (8,-0.2) {8};
\node[font=\sffamily\scriptsize\bfseries,text=tinta,anchor=south] at (4.50,2.05) {EM, iteración 8};
\end{tikzpicture} & \begin{tikzpicture}[x=0.42cm,y=0.62cm]
\draw[base] (0.4,0) -- (8.60,0);
\draw[base] (0.4,0) -- (0.4,2);
\draw[base] (0.4,0) -- (0.3,0); \node[font=\sffamily\tiny,text=gris,anchor=east] at (0.3,0) {0};
\draw[base] (0.4,1) -- (0.3,1); \node[font=\sffamily\tiny,text=gris,anchor=east] at (0.3,1) {1};
\draw[base] (0.4,2) -- (0.3,2); \node[font=\sffamily\tiny,text=gris,anchor=east] at (0.3,2) {2};
\node[anchor=south west,inner sep=0pt,text=nucC] at (0.560,0.0000) {\resizebox{0.370cm}{0.0196cm}{\sffamily\bfseries C}};
\node[anchor=south west,inner sep=0pt,text=nucG] at (0.560,0.0316) {\resizebox{0.370cm}{0.0411cm}{\sffamily\bfseries G}};
\node[anchor=south west,inner sep=0pt,text=nucT] at (0.560,0.0980) {\resizebox{0.370cm}{0.0416cm}{\sffamily\bfseries T}};
\node[anchor=south west,inner sep=0pt,text=nucA] at (0.560,0.1651) {\resizebox{0.370cm}{0.1471cm}{\sffamily\bfseries A}};
\node[font=\sffamily\tiny,text=gris] at (1,-0.2) {1};
\node[anchor=south west,inner sep=0pt,text=nucA] at (1.560,0.0065) {\resizebox{0.370cm}{0.0286cm}{\sffamily\bfseries A}};
\node[anchor=south west,inner sep=0pt,text=nucC] at (1.560,0.0526) {\resizebox{0.370cm}{0.0521cm}{\sffamily\bfseries C}};
\node[anchor=south west,inner sep=0pt,text=nucT] at (1.560,0.1366) {\resizebox{0.370cm}{0.8052cm}{\sffamily\bfseries T}};
\node[font=\sffamily\tiny,text=gris] at (2,-0.2) {2};
\node[anchor=south west,inner sep=0pt,text=nucT] at (2.560,0.0046) {\resizebox{0.370cm}{0.0372cm}{\sffamily\bfseries T}};
\node[anchor=south west,inner sep=0pt,text=nucC] at (2.560,0.0645) {\resizebox{0.370cm}{0.0851cm}{\sffamily\bfseries C}};
\node[anchor=south west,inner sep=0pt,text=nucG] at (2.560,0.2017) {\resizebox{0.370cm}{0.5890cm}{\sffamily\bfseries G}};
\node[font=\sffamily\tiny,text=gris] at (3,-0.2) {3};
\node[anchor=south west,inner sep=0pt,text=nucC] at (3.560,0.0154) {\resizebox{0.370cm}{0.0421cm}{\sffamily\bfseries C}};
\node[anchor=south west,inner sep=0pt,text=nucA] at (3.560,0.0833) {\resizebox{0.370cm}{0.9844cm}{\sffamily\bfseries A}};
\node[font=\sffamily\tiny,text=gris] at (4,-0.2) {4};
\node[anchor=south west,inner sep=0pt,text=nucC] at (4.560,0.0068) {\resizebox{0.370cm}{0.2825cm}{\sffamily\bfseries C}};
\node[anchor=south west,inner sep=0pt,text=nucG] at (4.560,0.4624) {\resizebox{0.370cm}{0.2953cm}{\sffamily\bfseries G}};
\node[font=\sffamily\tiny,text=gris] at (5,-0.2) {5};
\node[anchor=south west,inner sep=0pt,text=nucA] at (5.560,0.0035) {\resizebox{0.370cm}{0.0612cm}{\sffamily\bfseries A}};
\node[anchor=south west,inner sep=0pt,text=nucG] at (5.560,0.1022) {\resizebox{0.370cm}{0.0760cm}{\sffamily\bfseries G}};
\node[anchor=south west,inner sep=0pt,text=nucT] at (5.560,0.2247) {\resizebox{0.370cm}{0.4553cm}{\sffamily\bfseries T}};
\node[font=\sffamily\tiny,text=gris] at (6,-0.2) {6};
\node[anchor=south west,inner sep=0pt,text=nucG] at (6.560,0.0050) {\resizebox{0.370cm}{0.0243cm}{\sffamily\bfseries G}};
\node[anchor=south west,inner sep=0pt,text=nucA] at (6.560,0.0442) {\resizebox{0.370cm}{0.1307cm}{\sffamily\bfseries A}};
\node[anchor=south west,inner sep=0pt,text=nucC] at (6.560,0.2550) {\resizebox{0.370cm}{0.4448cm}{\sffamily\bfseries C}};
\node[font=\sffamily\tiny,text=gris] at (7,-0.2) {7};
\node[anchor=south west,inner sep=0pt,text=nucC] at (7.560,0.0000) {\resizebox{0.370cm}{0.0247cm}{\sffamily\bfseries C}};
\node[anchor=south west,inner sep=0pt,text=nucG] at (7.560,0.0398) {\resizebox{0.370cm}{0.0330cm}{\sffamily\bfseries G}};
\node[anchor=south west,inner sep=0pt,text=nucT] at (7.560,0.0930) {\resizebox{0.370cm}{0.0446cm}{\sffamily\bfseries T}};
\node[anchor=south west,inner sep=0pt,text=nucA] at (7.560,0.1649) {\resizebox{0.370cm}{0.6539cm}{\sffamily\bfseries A}};
\node[font=\sffamily\tiny,text=gris] at (8,-0.2) {8};
\node[font=\sffamily\scriptsize\bfseries,text=tinta,anchor=south] at (4.50,2.05) {EM, convergido};
\end{tikzpicture} & \begin{tikzpicture}[x=0.42cm,y=0.62cm]
\draw[base] (0.4,0) -- (8.60,0);
\draw[base] (0.4,0) -- (0.4,2);
\draw[base] (0.4,0) -- (0.3,0); \node[font=\sffamily\tiny,text=gris,anchor=east] at (0.3,0) {0};
\draw[base] (0.4,1) -- (0.3,1); \node[font=\sffamily\tiny,text=gris,anchor=east] at (0.3,1) {1};
\draw[base] (0.4,2) -- (0.3,2); \node[font=\sffamily\tiny,text=gris,anchor=east] at (0.3,2) {2};
\node[anchor=south west,inner sep=0pt,text=nucC] at (0.560,0.0000) {\resizebox{0.370cm}{0.0263cm}{\sffamily\bfseries C}};
\node[anchor=south west,inner sep=0pt,text=nucT] at (0.560,0.0424) {\resizebox{0.370cm}{0.0263cm}{\sffamily\bfseries T}};
\node[anchor=south west,inner sep=0pt,text=nucG] at (0.560,0.0849) {\resizebox{0.370cm}{0.0325cm}{\sffamily\bfseries G}};
\node[anchor=south west,inner sep=0pt,text=nucA] at (0.560,0.1373) {\resizebox{0.370cm}{0.1068cm}{\sffamily\bfseries A}};
\node[font=\sffamily\tiny,text=gris] at (1,-0.2) {1};
\node[anchor=south west,inner sep=0pt,text=nucC] at (1.560,0.0000) {\resizebox{0.370cm}{0.0072cm}{\sffamily\bfseries C}};
\node[anchor=south west,inner sep=0pt,text=nucA] at (1.560,0.0117) {\resizebox{0.370cm}{0.0362cm}{\sffamily\bfseries A}};
\node[anchor=south west,inner sep=0pt,text=nucG] at (1.560,0.0701) {\resizebox{0.370cm}{0.0362cm}{\sffamily\bfseries G}};
\node[anchor=south west,inner sep=0pt,text=nucT] at (1.560,0.1285) {\resizebox{0.370cm}{0.8182cm}{\sffamily\bfseries T}};
\node[font=\sffamily\tiny,text=gris] at (2,-0.2) {2};
\node[anchor=south west,inner sep=0pt,text=nucC] at (2.560,0.0067) {\resizebox{0.370cm}{0.0537cm}{\sffamily\bfseries C}};
\node[anchor=south west,inner sep=0pt,text=nucT] at (2.560,0.0932) {\resizebox{0.370cm}{0.0867cm}{\sffamily\bfseries T}};
\node[anchor=south west,inner sep=0pt,text=nucG] at (2.560,0.2330) {\resizebox{0.370cm}{0.3674cm}{\sffamily\bfseries G}};
\node[font=\sffamily\tiny,text=gris] at (3,-0.2) {3};
\node[anchor=south west,inner sep=0pt,text=nucG] at (3.560,0.0000) {\resizebox{0.370cm}{0.0075cm}{\sffamily\bfseries G}};
\node[anchor=south west,inner sep=0pt,text=nucT] at (3.560,0.0120) {\resizebox{0.370cm}{0.0075cm}{\sffamily\bfseries T}};
\node[anchor=south west,inner sep=0pt,text=nucC] at (3.560,0.0240) {\resizebox{0.370cm}{0.0671cm}{\sffamily\bfseries C}};
\node[anchor=south west,inner sep=0pt,text=nucA] at (3.560,0.1323) {\resizebox{0.370cm}{0.8424cm}{\sffamily\bfseries A}};
\node[font=\sffamily\tiny,text=gris] at (4,-0.2) {4};
\node[anchor=south west,inner sep=0pt,text=nucC] at (4.560,0.0142) {\resizebox{0.370cm}{0.2687cm}{\sffamily\bfseries C}};
\node[anchor=south west,inner sep=0pt,text=nucG] at (4.560,0.4475) {\resizebox{0.370cm}{0.2687cm}{\sffamily\bfseries G}};
\node[font=\sffamily\tiny,text=gris] at (5,-0.2) {5};
\node[anchor=south west,inner sep=0pt,text=nucG] at (5.560,0.0072) {\resizebox{0.370cm}{0.0584cm}{\sffamily\bfseries G}};
\node[anchor=south west,inner sep=0pt,text=nucA] at (5.560,0.1014) {\resizebox{0.370cm}{0.0764cm}{\sffamily\bfseries A}};
\node[anchor=south west,inner sep=0pt,text=nucT] at (5.560,0.2246) {\resizebox{0.370cm}{0.4178cm}{\sffamily\bfseries T}};
\node[font=\sffamily\tiny,text=gris] at (6,-0.2) {6};
\node[anchor=south west,inner sep=0pt,text=nucG] at (6.560,0.0074) {\resizebox{0.370cm}{0.0231cm}{\sffamily\bfseries G}};
\node[anchor=south west,inner sep=0pt,text=nucA] at (6.560,0.0447) {\resizebox{0.370cm}{0.1339cm}{\sffamily\bfseries A}};
\node[anchor=south west,inner sep=0pt,text=nucC] at (6.560,0.2607) {\resizebox{0.370cm}{0.4109cm}{\sffamily\bfseries C}};
\node[font=\sffamily\tiny,text=gris] at (7,-0.2) {7};
\node[anchor=south west,inner sep=0pt,text=nucC] at (7.560,0.0000) {\resizebox{0.370cm}{0.0069cm}{\sffamily\bfseries C}};
\node[anchor=south west,inner sep=0pt,text=nucT] at (7.560,0.0112) {\resizebox{0.370cm}{0.0069cm}{\sffamily\bfseries T}};
\node[anchor=south west,inner sep=0pt,text=nucG] at (7.560,0.0223) {\resizebox{0.370cm}{0.0899cm}{\sffamily\bfseries G}};
\node[anchor=south west,inner sep=0pt,text=nucA] at (7.560,0.1673) {\resizebox{0.370cm}{0.7538cm}{\sffamily\bfseries A}};
\node[font=\sffamily\tiny,text=gris] at (8,-0.2) {8};
\node[font=\sffamily\scriptsize\bfseries,text=tinta,anchor=south] at (4.50,2.05) {Gibbs, barrido 60};
\end{tikzpicture}
\end{tabular}
\caption[Un motivo emerge del ruido]{Un motivo emerge del ruido. \textbf{Arriba, izquierda}: logo de la matriz usada para implantar los sitios (8,11 bits en total; la columna 5 es una S = C/G). Los otros cinco logos siguen la corrida ganadora de EM desde una semilla aleatoria (\secuencia{TTCCCTTA}): la semilla apenas tiene 1,66 bits; en la iteración 6 empiezan a asomar la \nuc{G}, la \nuc{A} y la \nuc{T} centrales (3,19 bits); en la iteración 8 el motivo ya es reconocible (6,39 bits), y al converger (8,75 bits) coincide con el implantado. \textbf{Abajo, derecha}: la mejor cadena de Gibbs tras 60 barridos llega al mismo motivo por un camino estocástico. Alturas en bits según la ecuación del logo del \cref{cap:04}.}
\label{fig:13-logos}
\end{figure}

\subsection{Muestreo de Gibbs}

\textcite{c13-lawrence1993} abordaron el mismo problema con una estrategia estocástica, el \emph{muestreador de sitios} de Gibbs. En lugar de probabilidades fraccionarias, mantiene una asignación concreta $a=(a_1,\dots,a_N)$ de la posición del sitio en cada secuencia y la actualiza de una en una:

\begin{enumerate}
  \item Elija una secuencia $z$ y retire su sitio del alineamiento.
  \item Construya la matriz con los $N-1$ sitios restantes y seudoconteos:
  \begin{equation}\label{eq:13-gibbs-q}
  q_{k,b} = \frac{\sum_{i\ne z}\mathbb 1\bigl[x_{i,a_i+k-1}=b\bigr]+\beta_b}{N-1+\sum_{b'}\beta_{b'}} .
  \end{equation}
  \item Muestree una nueva posición para $z$ de la distribución condicional
  \begin{equation}\label{eq:13-gibbs}
  \Prob\bigl(a_z=j\mid a_{-z},X\bigr) \;\propto\; \prod_{k=1}^{W}\frac{q_{k,\,x_{z,j+k-1}}}{p_0(x_{z,j+k-1})},\qquad j=1,\dots,m .
  \end{equation}
  \item Repita para todas las secuencias (un \emph{barrido}) hasta que la verosimilitud se estabilice.
\end{enumerate}

\begin{simbolos}
\simbolo{$a_i$}{Posición asignada al sitio de la secuencia $i$.}
\simbolo{$a_{-z}$}{Posiciones de todas las secuencias excepto $z$.}
\simbolo{$q_{k,b}$}{Matriz estimada sin la secuencia $z$.}
\end{simbolos}

La diferencia con EM es de fondo. EM es determinista y sube siempre por la pendiente más cercana; Gibbs es una cadena de Márkov cuya distribución estacionaria es la posterior de las posiciones, así que puede \emph{bajar} momentáneamente y saltar a otra colina. En la práctica, las dos estrategias tienen los mismos enemigos: los óptimos locales y, sobre todo, los \emph{desfases} (\ingles{phase shifts}), en los que el alineamiento captura el motivo corrido una o dos posiciones. Los dos desfases de la \cref{fig:13-em}\,(b) (\secuencia{ACATGAGT} y \secuencia{CCATGAGT}) contienen seis columnas buenas de ocho y son muy estables, porque mover un único sitio no mejora nada; hace falta desplazar todos a la vez. \textcite{c13-lawrence1993} propusieron por ello movimientos que desplazan todo el alineamiento, y la práctica habitual es lanzar muchas cadenas y quedarse con la mejor. En nuestra simulación, la cadena ganadora alcanzó el \SI{95}{\percent} de su verosimilitud final en el sexto barrido y colocó 22 de los 30 sitios en la posición exacta.


\subsection{De motivos nuevos a factores conocidos: bases de datos y enriquecimiento}

Un motivo descubierto \emph{de novo} es solo una matriz. Para ponerle nombre se compara con bases de datos de perfiles de unión. JASPAR, en su novena versión, es la colección abierta de referencia de perfiles de factores de transcripción curados y no redundantes \parencite{c13-castromondragon2022}; la comparación entre matrices (por ejemplo, con la herramienta Tomtom de la MEME Suite, que puntúa la similitud columna a columna considerando desplazamientos y la hebra complementaria) devuelve los factores conocidos más parecidos con un valor de significancia \parencite{c13-bailey2009}.

\paragraph{Enriquecimiento diferencial} Una pregunta distinta, y a menudo más útil, es cuál de los miles de motivos conocidos aparece \emph{más} en los picos que en secuencias de fondo comparables. HOMER \parencite{c13-heinz2010} popularizó este enfoque con un fondo de secuencias genómicas aleatorias emparejadas por contenido de GC con los picos, lo que evita que cualquier motivo rico en GC parezca enriquecido solo porque los promotores lo son. Con ese enfoque, \textcite{c13-heinz2010} mostraron que el factor PU.1 colabora con pequeños conjuntos de factores determinantes de linaje para establecer los potenciadores propios de macrófagos y linfocitos B. La significancia se calcula con la distribución hipergeométrica:
\begin{equation}\label{eq:13-hiper}
p = \Prob(K\ge k_t) = \sum_{k\ge k_t}\frac{\binom{k_t+k_b}{k}\binom{n_t+n_b-k_t-k_b}{n_t-k}}{\binom{n_t+n_b}{n_t}},
\end{equation}
\begin{simbolos}
\simbolo{$n_t,\ k_t$}{Número de secuencias objetivo (picos) y cuántas contienen el motivo.}
\simbolo{$n_b,\ k_b$}{Número de secuencias de fondo y cuántas lo contienen.}
\simbolo{$K$}{Número de secuencias con motivo que caerían en el conjunto objetivo si la asignación fuera al azar.}
\end{simbolos}
\noindent es decir, la probabilidad de que, repartiendo al azar las $k_t+k_b$ secuencias con motivo entre objetivo y fondo, al menos $k_t$ caigan en el objetivo.

\begin{ejemplo}{Un motivo enriquecido}{13-homer}
De \num{2000} picos, 640 (\SI{32}{\percent}) contienen el motivo; de \num{20000} secuencias de fondo emparejadas por GC, lo contienen \num{1200} (\SI{6}{\percent}). El enriquecimiento es $0{,}32/0{,}06=5{,}33$ veces, y la \cref{eq:13-hiper} da $p=6{,}8\times10^{-234}$ ($\ln p=-536{,}9$, la escala que HOMER reporta). La aproximación binomial, que trata el fondo como una proporción conocida sin error, da un valor aún menor ($2\times10^{-276}$). Valores tan extremos son habituales y no deben leerse literalmente: con miles de secuencias, cualquier sesgo sistemático de composición da $p$ astronómicos. Lo que importa es el \emph{orden} de los motivos y el tamaño del enriquecimiento.
\end{ejemplo}

\subsection{Centralidad: ¿unión directa o indirecta?}

Un motivo enriquecido en los picos no es necesariamente el del factor inmunoprecipitado: puede ser el de un cofactor que se une cerca, o reflejar un sesgo de las regiones abiertas. \textcite{c13-bailey2012} observaron que, si el factor se une \emph{directamente} al ADN, sus sitios deben concentrarse en el centro de los picos (la posición del \ingles{summit}), mientras que los de cofactores se distribuyen más anchos y los del fondo, uniformemente. Su herramienta CentriMo busca la ventana central donde el enriquecimiento es máximo y lo contrasta con una prueba binomial: si una ventana de ancho $w_c$ ocupa una fracción $f=w_c/(L-W+1)$ de las posiciones posibles, el número de mejores sitios que caen dentro sigue, bajo la hipótesis nula, una $\mathrm{Bin}(n,f)$. Así mostraron, por ejemplo, que en un conjunto de datos de Smad1 no había evidencia de unión directa del factor inmunoprecipitado.

\begin{figure}[t]
\centering
\begin{tikzpicture}
\begin{axis}[curso, width=10.5cm, height=5.4cm,
   xlabel={posición del mejor sitio respecto al centro del pico (pb)},
   ylabel={número de picos}, xmin=-250, xmax=250, ymin=0, ymax=820,
   xtick={-250,-150,-50,0,50,150,250},
   legend style={at={(0.01,0.97)},anchor=north west}]
  \fill[amarillo!20] (axis cs:-50,0) rectangle (axis cs:50,820);
  \addplot[azul, very thick, const plot mark mid] table[x=x,y=dir] {figuras/cap13/centrimo.dat};
  \addlegendentry{factor directo (97,6\,\% en $\pm50$)}
  \addplot[naranja, very thick, const plot mark mid] table[x=x,y=cof] {figuras/cap13/centrimo.dat};
  \addlegendentry{cofactor (37,7\,\%)}
  \addplot[gris, very thick, const plot mark mid] table[x=x,y=fondo] {figuras/cap13/centrimo.dat};
  \addlegendentry{motivo sin relación (21,5\,\%)}
\end{axis}
\end{tikzpicture}
\caption[Análisis de centralidad]{Análisis de centralidad al estilo de CentriMo sobre \num{3000} picos simulados de 500 pb. Los mejores sitios del motivo del factor inmunoprecipitado (azul) se concentran en los 100 pb centrales (banda amarilla), donde cae el \SI{97,6}{\percent} de ellos frente al \SI{21}{\percent} esperado; la prueba binomial da un valor $p$ numéricamente nulo. El motivo de un cofactor (naranja) está enriquecido, pero de forma mucho más ancha ($p=2{,}9\times10^{-35}$), y un motivo sin relación se distribuye de forma plana ($p=0{,}39$). La centralidad distingue unión directa de indirecta, algo que el simple enriquecimiento no puede hacer.}
\label{fig:13-centrimo}
\end{figure}


\begin{consola}[Análisis de motivos a partir de picos]
# 100 pb alrededor de la cumbre de cada pico
awk -v OFS='\t' '{s=$2+$10-50; print $1,s,s+100}' \
    ctcf_idr.narrowPeak > cumbres.bed
bedtools getfasta -fi hg38.fa -bed cumbres.bed -fo cumbres.fa
# MEME Suite: de novo + enriquecimiento + centralidad
meme-chip -oc memechip -meme-nmotifs 5 \
    -db JASPAR2022_CORE_vertebrates.meme cumbres.fa
# HOMER: fondo emparejado por GC
findMotifsGenome.pl cumbres.bed hg38 homer_out -size given
\end{consola}

\begin{cuidado}[Motivos que engañan]
\begin{enumerate}
  \item \emph{Baja complejidad y repeticiones.} Tramos como \secuencia{CACACACA} o poli-A se ``descubren'' con enorme significancia; enmascare repeticiones y use modelos de fondo de orden superior.
  \item \emph{Sesgo de GC.} Los promotores son ricos en GC; sin un fondo emparejado, cualquier motivo rico en GC parecerá enriquecido.
  \item \emph{Presencia no es ocupación.} Un genoma humano contiene cientos de miles de coincidencias de casi cualquier PWM; la inmensa mayoría no está ocupada, porque la accesibilidad de la cromatina (\cref{sec:13-chip}) decide.
\end{enumerate}
\end{cuidado}

\begin{ideaclave}
Descubrir motivos es estimar una PWM con datos incompletos: EM reparte votos fraccionarios y sube por la verosimilitud; Gibbs muestrea posiciones y puede escapar de óptimos locales. Ninguno garantiza el óptimo global: múltiples arranques, fondos emparejados y el criterio de centralidad convierten una matriz en una hipótesis biológica.
\end{ideaclave}

\begin{resumen}
\begin{itemize}
  \item En ChIP-seq, las lecturas de las hebras $+$ y $-$ flanquean el sitio a una distancia igual al tamaño de fragmento, que se estima por correlación cruzada; NSC y RSC comparan el pico de fragmento con el pico fantasma de mapabilidad.
  \item MACS llama picos con una prueba de Poisson cuya media es el máximo entre el fondo global y varias estimaciones locales del control; FRiP, la lista negra de ENCODE e IDR evalúan la señal y la reproducibilidad.
  \item ATAC-seq mide la accesibilidad con Tn5; un buen experimento muestra fragmentos libres de nucleosomas, jorobas nucleosomales, periodicidad helicoidal y un fuerte enriquecimiento en el TSS. CUT\&RUN reduce el fondo llevando la nucleasa al factor.
  \item El bisulfito convierte C no metiladas en T; Bismark alinea en un alfabeto reducido. La metilación es una proporción ($\beta$), su logit es el valor $M$, y la comparación entre grupos requiere un modelo beta-binomial que incluya la variación biológica.
  \item El descubrimiento de motivos se formula como un problema de datos incompletos: EM (MEME) y muestreo de Gibbs aprenden a la vez la PWM y las posiciones; el enriquecimiento diferencial (HOMER) y la centralidad (CentriMo) identifican los factores responsables.
\end{itemize}
\end{resumen}

\begin{ejercicios}
\ej{1} Una lectura de 50 pb de la hebra $-$ tiene su extremo 5' en la coordenada \num{10480}. Si el tamaño de fragmento estimado es $d=180$, ¿en qué coordenada colocaría el centro del fragmento? ¿Y si la lectura fuera de la hebra $+$ con extremo 5' en \num{10480}?
\ej{1} Calcule $\lambda_{\text{BG}}$ para un experimento de \num{3e7} lecturas en ratón ($G=\num{1.87e9}$) con $d=150$. ¿Cuántas lecturas en una ventana hacen falta para $p<10^{-5}$? Compruébelo con \py{scipy.stats.poisson.sf}.
\ej{2} Demuestre que, si $M=\log_2\frac{\beta}{1-\beta}$, entonces $\beta=\frac{2^M}{1+2^M}$, y use el método delta para mostrar que $\operatorname{Var}(\hat M)$ es mínima en $\beta=0{,}5$. ¿Qué valor mínimo toma con $n=30$ lecturas?
\ej{2} Simule \num{10000} CpG con $\mu=0{,}7$, tres réplicas por grupo y cobertura 20, bajo la hipótesis nula (ambos grupos iguales) con $\phi=0{,}1$. Aplique la prueba de Wald con $\phi=0$ y con $\phi=0{,}1$. ¿Qué fracción de valores $p$ cae por debajo de 0,05 en cada caso? Relacione el resultado con la \cref{eq:13-betabin}.
\ej{2} Modifique \py{em_oops} para el modelo ZOOPS añadiendo, para cada secuencia, un estado ``sin sitio'' con probabilidad a priori $1-\gamma$. Escriba el paso E y el paso M para $\gamma$.
\ej{3} Implemente el muestreador de Gibbs de las \cref{eq:13-gibbs-q,eq:13-gibbs} y añada un movimiento de desfase que proponga desplazar todas las posiciones $\pm1$ y lo acepte según la razón de verosimilitudes. Compare, con 20 cadenas, la fracción que alcanza el motivo correcto con y sin ese movimiento.
\ej{3} Demuestre la desigualdad que sustenta la monotonía de EM: para cualquier distribución $q(Z)$, $\log\Prob(X\mid\theta)=\mathcal L(q,\theta)+\mathrm{KL}\bigl(q\,\|\,\Prob(Z\mid X,\theta)\bigr)$, donde $\mathcal L$ es la cota de Jensen. Deduzca que el paso E anula la divergencia y que el paso M no puede disminuir la verosimilitud.
\ej{3} Descargue de ENCODE los picos IDR de CTCF en una línea celular, extraiga $\pm50$ pb alrededor de cada cumbre y ejecute \cmd{meme-chip}. ¿Coincide el motivo principal con el perfil MA0139 de JASPAR? ¿Qué fracción de los picos lo contiene y qué muestra el análisis de centralidad?
\end{ejercicios}

\referenciascapitulo
